\section{Introduction and exact scope}

Write
\[
 \Ein(z)=\int_0^z\frac{1-e^{-t}}t\,dt
 =\sum_{r\ge1}\frac{(-1)^{r-1}z^r}{r\,r!},
\]
and let
\[
 \gamma=\lim_{n\to\infty}\left(\sum_{k=1}^n\frac1k-\log n\right),
 \qquad
 \delta=eE_1(1).
\]
Neither \(\gamma\) nor \(\delta\) is known to be irrational; see
Lagarias' survey \cite{Lagarias2013}.  The elementary positive-axis identity
\begin{equation}\label{eq:Ein-connection}
 \Ein(x)=\gamma+\log x+E_1(x)\qquad(x>0)
\end{equation}
places \(\gamma\) at the boundary between algebraic-point \(E\)-values and an
asymptotic connection constant.

The companion paper \cite{KEIO2026I} determined the algebraic relations among
all nonzero algebraic-point \(\Ein\)-values over the pure exponential field,
constructed several finite Euler witnesses, and computed the first two
inverse traces of the divisor of \(\Ein-\gamma\).  The present paper begins
where that finite theory stops.  Its organizing question is:
\begin{quote}
What remains true when the level of the divisor is allowed to move through a
countable algebraically independent family converging to \(\gamma\)?
\end{quote}

The results have four layers.  In the spectral layer, the two even traces
of orders two and four
recover the level exactly, so the entire fixed-level inverse-trace ring is
one-dimensional; moving the level restores one algebraically
independent coordinate per level.  The order-two row has a positive
telescoping structure and hence a genuine operator-theoretic realization.
Algebraic levels of a natural gamma-free inverse map have
transcendental preimages.  All finite linear witnesses are governed by
one normalized-certificate uniqueness lemma.  The dyadic tail is the
unique minimizer in a natural cumulative-positive augmentation cone, while
outside that cone three-point forms are already dense.  The resulting
coefficient-height transition identifies the exact scale at which the first
tail can be cancelled.  In the transmutation layer, the four-term Euler
recurrence from the 2026 Ramanujan Challenge is connected explicitly to the
official Meijer--\(G\) construction, the cubic Pilehrood family, a finite
probability law, a regularized Gamma determinant, and a full rational
zeta-jet tower.  Its denominator and numerator live in a flat dual-number
\(SL_3\) system, equivalently an ordinary \(SL_6\) system.  In the
obstruction layer, the particular residual-prime,
fixed-place, Rodrigues, natural-index P-recursive, and strict-Stokes
constructions studied here yield scoped saturation identities, budget
inequalities, or sharply stated missing inputs.

The principal results are as follows.

\begin{theorem}[Universal trace and transversal independence]
\label{thm:intro-universal}
For \(c\in\C^\times\), let \(Z_c\) be the zero multiset of
\(\Ein(z)-c\).  For every \(m\ge2\),
\[
 T_m(c):=\sum_{\rho\in Z_c}\rho^{-m}=P_m(c^{-1}),
\]
where \(P_m\in\Q[X]\) is monic of degree \(m\).  If
\(\{c_j:j\in J\}\) is algebraically independent over a characteristic-zero
field \(F\), and one integer \(m_j\ge2\) is selected for each \(j\), then
\(\{T_{m_j}(c_j):j\in J\}\) is algebraically independent over \(F\).
Moreover
\[
 c^{-1}=144T_4(c)-144T_2(c)^2-14T_2(c),
\]
so \(F[T_m(c):m\ge2]=F[c^{-1}]\).
The map \(c\mapsto(T_2(c),T_3(c))\) is injective except for the single
unordered pair
\[
 \{6-2\sqrt3,6+2\sqrt3\},
\]
whose common trace pair is \((-1/24,1/96)\).  For each \(m\ge2\), the
nonzero roots of \(P_m\) also give a nonempty finite set of algebraic levels
at which \(T_m\) vanishes and every point of the corresponding divisor is
transcendental.
\end{theorem}

\begin{theorem}[Finite-rank tails and incomplete-Gamma jets]
\label{thm:intro-tail-closures}
Let \(\Gamma\subset\Qbar_{>0}^{\times}\) have multiplicative rank \(r\).
For every finite \(A\subset\Gamma\),
\[
 \trdeg_KK(E_1(\alpha):\alpha\in A)\ge |A|-r-1.
\]
After deleting at most \(r+1\) values, the remaining full \(E_1\)-tail
family is jointly algebraically independent over \(K\).  The same bound
holds for the combined classical family
\(\{E_1(\alpha),\Ei(\alpha),\Ci(\alpha),\operatorname{Si}(\alpha)\}\).

For the regularized lower incomplete-Gamma jets
\[
 R_{\alpha,k}=[s^k]\{s^{-1}-\Gamma_{<}(s,\alpha)\},
\]
every finite \(S\subset\Gamma\times\N_0\) satisfies
\[
 \trdeg_KK(R_{\alpha,k}:(\alpha,k)\in S)\ge |S|-r.
\]
Thus at most \(r\) exceptions need be removed from the complete two-index
lower-jet family.  The corresponding upper jets have defect at most
\(r\) plus the number of distinct jet orders.
\end{theorem}

\begin{theorem}[Dyadic moving zero-trace tower]
\label{thm:intro-dyadic}
For \(N\ge1\), put
\[
 Y_N=(N+1)\Ein(2^N)-N\Ein(2^{N+1}).
\]
Then \(\{Y_N:N\ge1\}\) is countably algebraically independent over \(K\),
\(Y_N\downarrow\gamma\), and
\[
 Y_N-\gamma\sim \frac{N+1}{2^N}e^{-2^N}.
\]
For every choice \(m_N\ge2\), the exact traces
\[
 \tau_{m_N,N}:=\sum_{\Ein(\rho)=Y_N}\rho^{-m_N}
\]
are countably algebraically independent over \(K\).  For fixed \(m\),
\(\tau_{m,N}\to S_m\), the \(m\)-th inverse trace of \(\Ein-\gamma\).
\end{theorem}

\begin{theorem}[Inverse levels and critical coefficients]
\label{thm:intro-inverse}
The function
\[
 \mathscr R(q)=E_1(1)-2E_1(q)+E_1(q^2)
\]
maps \((1,\infty)\) increasingly onto \((0,E_1(1))\).  The inverse image of
every algebraic number in this interval is transcendental.  In particular,
the unique roots \(r_m>1\) of \(\mathscr R(r_m)=1/m\), \(m\ge5\), are all
transcendental.

For each fixed algebraic \(q>1\), the critical three-point cancellation
coefficients \(u_N(q)\) defined in \eqref{eq:uNq} satisfy, for every finite
\(I\),
\[
 \trdeg_KK(u_N(q):N\in I)\ge |I|-1.
\]
\end{theorem}

\begin{theorem}[Positive spectral realizations]
\label{thm:intro-positive}
Set
\[
 \tau_N=\frac1{Y_N^2}-\frac1{2Y_N},\qquad
 d_0=\tau_1,\qquad d_N=\tau_{N+1}-\tau_N\quad(N\ge1).
\]
Then \(d_N>0\), the \(d_N\) are countably algebraically independent over
\(K\), and
\[
 S_2=\sum_{N\ge0}d_N.
\]
The diagonal operator \(\mathsf A e_N=d_Ne_N\) belongs to every Schatten
class and has trace \(S_2\).  Moreover
\[
 \Delta_2(z)=\det\!\left(I-\frac{z^2}{2}\mathsf A\right)
 =\prod_{N\ge0}\left(1-\frac{d_N}{2}z^2\right)
\]
is an order-zero Laguerre--P\'olya entire function.  Its positive zeros are
countably algebraically independent over \(K\), and
\[
 \sum_{\rho\in Z(\Delta_2)}\rho^{-2}=S_2.
\]
The analogous order-four increments define a positive diagonal operator
\(\mathsf A_4\) with trace \(S_4\), and, writing
\(\mathsf A_2=\mathsf A\),
\[
 \gamma=
 \left\{144\Tr\mathsf A_4
 -144(\Tr\mathsf A_2)^2-14\Tr\mathsf A_2\right\}^{-1}.
\]
\end{theorem}

\begin{theorem}[Cross-index witness rigidity]
\label{thm:intro-witness}
Let \(q_*\in(0,1)\) be the real solution of \(\Ein(q_*)=\gamma\).  Let
\(\beta_n\in(0,1)\) be the positive real zero of
\[
 G_n(z)=n\Ein(z)-\Ein(z^n)-(n-1)\gamma,
\]
and, for odd \(n\ge3\), let \(\lambda_n<-1\) be its negative real zero.
Then at most one member of
\[
 \{q_*\}\cup\{\beta_n:n\ge2\}
 \cup\{\lambda_n:n\ge3,\ n\text{ odd}\}
\]
is algebraic.
\end{theorem}

\begin{theorem}[Critical moving-point barrier]
\label{thm:intro-barrier}
Any theorem of the natural form
\[
 |P(F(\xi))|\ge
 \exp\!\bigl(-C(1+|\xi|)^c\Phi(\deg P,\log H(P))\bigr)
\]
intended to hold uniformly for all fixed \(E\)-function systems and positive
integral sample points, or already for all nonzero polynomial values of the
displayed system, must have \(c\ge1\).  The obstruction occurs for the fixed
system
\[
 F=(\Ein(z),\Ein(2z),\Ein(4z))
\]
and the fixed linear form \(P=X_1-2X_2+X_3\).  General sufficiency at
\(c=1\) is not known; at that scale the same example forces at least the
polynomial loss \(\xi^{-1}\).  Within the augmentation-one stratum the
dyadic coefficients uniquely minimize the tail in the cumulative-positive
cone.  Outside that cone three points give dense tail values, while
cancellation of the first active tail requires at least
\(\log H\ge(q-1)q^N+O(1)\).  For consecutive augmentation-one triples this
transition has a matching construction; at natural point cost a normalized
lower bound for that affine template must pay height exponent at least two.

For the dyadic tower, algebraicity of \(S_2\) is
equivalent to the existence of one algebraic coefficient \(s\) for which
\[
 -\log|2sY_N^2+Y_N-2|\sim2^N.
\]
\end{theorem}

\begin{theorem}[Exact Ramanujan transmutation package]
\label{thm:intro-recurrence}
The two distinguished solutions of the Euler recurrence
\eqref{eq:challenge-recurrence} have the finite sums
\eqref{eq:qn-closed}--\eqref{eq:pn-closed}, satisfy the recurrence for every
\(n\ge0\), and obey \(p_n/q_n\to\gamma\).  Their terminating
\({}_2F_2\) state and the official Meijer--\(G\) state are joined by the
rational coboundary \eqref{eq:ram-coboundary}; the normalized conserved
pairing is the finite identity \eqref{eq:ram-bilinear-matrix}.  The
denominator system embeds in the flat \(SL_3\) field
\eqref{eq:Ma}--\eqref{eq:cmf-flatness}, and its zero--first jet lift is the
flat \(SL_6\) system \eqref{eq:SL6-block}.

The Euler--Beta--log transform
\eqref{eq:q-Beta-transform}--\eqref{eq:p-Beta-transform} identifies the
same solutions with the cubic Pilehrood family.  The rational carrier
\(\mathscr K_a\) has denominator and numerator as its zero- and first-jets;
\(\mathscr K_a(s)/\mathscr K_a(0)\) converges to
\(\Gamma(1-s)^3\Gamma(1+s)^2\) and yields the zeta limits
\eqref{eq:zeta-log-jets}.  Its scaled integer numerator is the monic contact
polynomial of \cref{thm:contact-polynomial}.  All recurrence solutions have
the limiting form \eqref{eq:all-initial-limits}.  Moreover, with
\(b=n+4\) and \(D_b=\lcm(1,\ldots,b)\),
\[
 \frac{b!}{D_b}\gcd\!\left(b,\frac{D_b}{b}\right)
 \mid\gcd(p_n,q_n),
\]
and the factorial-six part of the exact Casoratian
\eqref{eq:challenge-casoratian} comes from the determinant of the
denominator-cleared official dual transfer.
\end{theorem}

\begin{theorem}[Arithmetic and Stokes barriers]
\label{thm:intro-new-barriers}
The following statements hold unconditionally.
\begin{enumerate}[label=\textup{(\roman*)}]
 \item The dyadic level \((Y_N)\) and its order-two trace \((\tau_N)\) are
 not P-recursive.
 \item The two-base form of \cref{prop:two-base-obstruction} remains
 exponentially small with
 \[
  -\log\{Y_N(2)-Y_{M_N}(3)\}
  =2^N+N\log2-\log(N+1)+O(1).
 \]
 Thus multiplicative independence of the two bases does not remove this
 moving-point scale.
 \item Every finite multi-level trace field is exactly the
 \(\mathbb G_a\)-invariant field of the gauge shift
 \(G\mapsto G+t,H_i\mapsto H_i-t\).  The completed moving tower breaks this
 formal symmetry through its boundary limit.
 \item Every nondegenerate rank-one real Hilbert quotient obtained from the
 fixed residual-direction constraints is exactly saturated under a candidate
 \(\gamma=a/b\).  For the optimally saddled uncombined Pr\'evost form, the
 lcm-only height-one ledger has deficit
 \(2-\log(3+2\sqrt2)\).
 \item The fixed-prime Rodrigues family \eqref{eq:fixed-p-rodrigues}
 places \(p^{n-\ell}\) in the endpoint coefficient.  Even if this whole
 factor is credited as fixed-place gain, the resulting ledger and the
 archimedean saddle satisfy \eqref{eq:universal-fixed-p-bound}.
 \item The displayed unframed lateral transition of the elementary
 differential operator annihilating \(F=-\Ein\) is gamma-free.  The Euler
 constant occurs in the connection framing and, under
 \(\gamma\in\Qbar\), can be removed by an algebraic shear.  A framed
 Euler--Tate matched-line lifting theorem would imply the
 \(\Qbar\)-linear independence of \(1,\pi,\gamma\), and hence the
 transcendence of \(\gamma\), \(\gamma/\pi\), and \(S_2\).
\end{enumerate}
\end{theorem}

The qualification ``transversal'' in
\cref{thm:intro-universal,thm:intro-dyadic} is essential: one chooses one
trace order at each level.  At a fixed level all trace orders are algebraic
over the same one-dimensional level field.  Likewise, the infinite sum in
\cref{thm:intro-positive} has no automatic arithmetic consequence.  An
infinite sum of algebraically independent transcendental numbers may be
algebraic.

\paragraph{Organization.}
\Cref{sec:input} records the exact input from the companion paper.
\Cref{sec:rank-deficit} proves finite-rank cofinite independence for both
connection coordinates and classical tails, together with the complete
cancellation geometry; \cref{sec:external-realizations} exports the
engine to sine-integral graphs, generalized Dickman transforms, and
regularized incomplete-Gamma jets; the
characteristic-zero algebraic-matroid realization is isolated in
\cref{app:Ein-matroids} because it uses an external linearization theorem.
\Cref{sec:trace-polynomials} develops the universal trace calculus and
identifies its finite multi-level field with the exact shift quotient.
\Cref{sec:dyadic} proves the dyadic and all-order moving towers.
\Cref{sec:spectral} constructs the two positive operators and the
order-two Laguerre--P\'olya determinant.  \Cref{sec:other-towers} gives fixed-coefficient
power-orbit towers, transcendental inverse levels, and the logarithmic
lattice.  \Cref{sec:certificates} proves normalized certificate uniqueness,
cross-index rigidity, and the one-defect cancellation theorem.
\Cref{sec:relative} gives the relative five-class exclusion and the exact
finite-descent criterion.
\Cref{sec:ramanujan-recurrence,sec:cmf} solve the four-term Euler recurrence,
classify all of its limits, and construct the flat \(SL_3/SL_6\) field.
\Cref{sec:meijer-duality,sec:beta-log} give the exact official gauge and the
Euler--Beta--log offset ladder.  \Cref{sec:gamma-jets,sec:ram-probability}
develop the rational Gamma jets, contact polynomial, finite probability
law, and zero geometry; \cref{sec:gamma-determinant} gives the regularized
determinant realization.  \Cref{sec:cross-place-pade} proves the real and
\(p\)-adic Euler--Pad\'e convergence results.
\Cref{sec:barrier} proves the sharp inside/outside augmentation dichotomy and
isolates the critical moving-point obstruction.
\Cref{sec:sondow} replaces an invalid asymmetric-integral heuristic by an
exact evaluation.  \Cref{sec:prevost-grid,sec:fixed-prime-rodrigues}
establish the residual-prime, Ridout, and Rodrigues barriers.
\Cref{sec:stokes-framing,sec:mixed-gevrey,sec:euler-tate} separate unframed
Stokes data from the mixed arithmetic-Gevrey and canonical motivic
framings, prove the conditional dilation and Gamma-jet ranks, and isolate
the strengthened matched-line target.  The final
section records what remains open.

\paragraph{Dependency and status ledger.}
To keep imported input, internal proof, and conditional extrapolation
separate, the logical dependencies are summarized below.
\begin{center}
\small
\begin{tabular}{@{}
 >{\raggedright\arraybackslash}p{0.24\textwidth}
 >{\raggedright\arraybackslash}p{0.32\textwidth}
 >{\raggedright\arraybackslash}p{0.34\textwidth}@{}}
\toprule
Source & Imported statement & Status in this paper\\
\midrule
Companion paper \cite[Thms.~1.4, 5.1 and Lem.~11.2]{KEIO2026I}
& Algebraic independence of algebraic-point resonant \(\Ein_q\)-values
and of the balanced dyadic coordinates
& Standing arithmetic input; the tail, incomplete-Gamma, trace, carrier,
and rigidity deductions are proved here\\
Fischler--Rivoal \cite{FischlerRivoalMixed2024}
& Definitions of mixed functions; conjecture
\(\mathbf E\cap\mathbf D=\Qbar\); conditional linear specialization
& The dilation and all-order Gamma-jet families are constructed here;
every conclusion using the intersection conjecture is explicitly
conditional\\
Official Challenge solution \cite[\S5.2]{RamanujanOfficial2026}
& Official Meijer--\(G\) normalization and Miller-selected limiting direction
& Imported normalization; the rational coboundary, bilinear identity, and
transmutations are proved here\\
Rosen--Sidman--Theran \cite[Thm.~1.10]{RosenSidmanTheran2025}
& Linear representability over an algebraically closed characteristic-zero
field
& Used only in Appendix~\ref{app:Ein-matroids}\\
Powers \cite[\S5]{Powers2025}
& A concluding multisection algebraic-independence assertion
& Used only under an explicit hypothesis in
Remark~\ref{rem:Powers-conditional}; no unconditional result depends on it\\
\bottomrule
\end{tabular}
\end{center}

\section{Input from the companion paper}
\label{sec:input}

Throughout,
\[
 K=\Eexp=\Qbar(e^\alpha:\alpha\in\Qbar),
\]
and \(K^{\mathrm{alg}}\) denotes its relative algebraic closure in \(\C\).
We use the following theorem from \cite[Thm.~5.1]{KEIO2026I}.

\begin{theorem}[Algebraic-point \(\Ein\)-independence]
\label{thm:input-independence}
The countable family
\[
 \{\Ein(\alpha):\alpha\in\Qbar^\times\}
\]
is algebraically independent over \(K\), in the finite-subset sense.
Consequently it remains algebraically independent over
\(K^{\mathrm{alg}}\).
\end{theorem}

For the incomplete-Gamma application below we also use the resonant
extension of this input.  Put
\[
 \Ein_q(z)=\sum_{n\ge1}\frac{(-1)^{n-1}z^n}{n^q n!}
 \qquad(q\ge1).
\]
By \cite[Thm.~1.4]{KEIO2026I}, the family
\begin{equation}\label{eq:resonant-input}
 \{\Ein_q(\alpha):q\ge1,\ \alpha\in\Qbar^\times\}
\end{equation}
is algebraically independent over \(K\), in the finite-subset sense.

The final assertion is standard: algebraic independence is preserved under
algebraic extension of the base field.  We shall repeatedly use the following
linear consequence.

\begin{lemma}[Full-rank linear images]
\label{lem:linear-images}
Let \(x_1,\ldots,x_r\) be algebraically independent over a field \(F\), and
let \(A\in M_{s\times r}(F)\) have row rank \(s\).  Then the \(s\) coordinates
of \(A(x_1,\ldots,x_r)^{\mathsf T}\) are algebraically independent over \(F\).
\end{lemma}

\begin{proof}
Complete the rows of \(A\) to an invertible \(r\)-by-\(r\) matrix over \(F\).
The resulting linear change of variables is an automorphism of the rational
function field \(F(x_1,\ldots,x_r)\).
\end{proof}

For later comparison we also recall the Euler-level traces from
\cite[Thm.~14.5]{KEIO2026I}.  If \(Z_\gamma\) is the zero multiset of
\(\Ein(z)-\gamma\), then
\begin{align}
 S_2&:=\sum_{\rho\in Z_\gamma}\rho^{-2}
 =\frac1{\gamma^2}-\frac1{2\gamma},\label{eq:S2}\\
 S_3&:=\sum_{\rho\in Z_\gamma}\rho^{-3}
 =\frac1{\gamma^3}-\frac3{4\gamma^2}+\frac1{6\gamma},\label{eq:S3}
\end{align}
and
\begin{equation}\label{eq:gamma-recovery}
 \gamma=\frac{24S_2+1}{24S_3+6S_2}.
\end{equation}
In particular \(S_2\) is algebraic if and only if \(\gamma\) is algebraic.

\section{Finite-rank independence and cancellation geometry}
\label{sec:rank-deficit}

The algebraic-point theorem has a useful export form when the arguments lie
in a multiplicative group of finite rank.  We restrict here to positive
algebraic arguments, so that the real logarithm is single-valued and no
additional \(2\pi i\)-coordinate occurs.  Put
\begin{equation}\label{eq:c-alpha}
 c_\alpha=\Ein(\alpha)-\log\alpha=\gamma+E_1(\alpha)
 \qquad(\alpha\in\Qbar_{>0}).
\end{equation}

\begin{theorem}[Finite-rank deficit and cofinite independence]
\label{thm:cofinite-independence}
Let \(\Gamma\subset\Qbar_{>0}^{\times}\) have multiplicative rank \(r\).
For every finite \(A\subset\Gamma\),
\begin{equation}\label{eq:finite-rank-deficit}
 \trdeg_KK(c_\alpha:\alpha\in A)\ge |A|-r.
\end{equation}
Consequently there is a set \(E\subset\Gamma\), \(|E|\le r\), such that
\begin{equation}\label{eq:cofinite-family}
 \{c_\alpha:\alpha\in\Gamma\setminus E\}
 \quad\hbox{is algebraically independent over }K.
\end{equation}

Fix \(\alpha_0\in\Gamma\).  There is likewise a set
\(E_{\alpha_0}\subset\Gamma\setminus\{\alpha_0\}\),
\(|E_{\alpha_0}|\le r\), for which
\begin{equation}\label{eq:cofinite-differences}
 \{E_1(\alpha)-E_1(\alpha_0):
 \alpha\in\Gamma\setminus(E_{\alpha_0}\cup\{\alpha_0\})\}
\end{equation}
is algebraically independent over \(K\).
\end{theorem}

\begin{proof}
Choose \(\lambda_1,\ldots,\lambda_r\) as a basis of the \(\Q\)-vector space
spanned by \(\{\log\alpha:\alpha\in\Gamma\}\); torsion causes no difficulty
because all elements are positive.  For every finite \(A\), all
\(\log\alpha\), \(\alpha\in A\), lie in the \(\Q\)-span of the
\(\lambda_j\).  By \eqref{eq:c-alpha},
\[
 K(\Ein(\alpha):\alpha\in A)
 \subset K(c_\alpha:\alpha\in A,\lambda_1,\ldots,\lambda_r).
\]
The left side has transcendence degree \(|A|\) by
\cref{thm:input-independence}; this proves \eqref{eq:finite-rank-deficit}.

Choose a maximal algebraically independent subfamily of
\(\{c_\alpha:\alpha\in\Gamma\}\).  If its complement contained
\(r+1\) elements, the algebraic dependence of each of them on a finite part
of the maximal family would produce a finite set with deficit at least
\(r+1\), contradicting \eqref{eq:finite-rank-deficit}.  This proves
\eqref{eq:cofinite-family}.

For distinct \(\alpha_1,\ldots,\alpha_m\ne\alpha_0\), the differences
\(\Ein(\alpha_i)-\Ein(\alpha_0)\) are algebraically independent over \(K\):
they are full-rank linear images of the independent coordinates at
\(\alpha_0,\alpha_1,\ldots,\alpha_m\).  Subtracting the logarithmic
differences costs at most the same \(r\) logarithmic coordinates.  The first
argument and the maximality argument therefore apply verbatim.
\end{proof}

The common Euler term costs at most one further transcendence degree.  This
turns the shift bookkeeping into an unconditional theorem about the tails
themselves.

\begin{theorem}[Finite-rank tail defect and cofinite independence]
\label{thm:E1-cofinite}
Let \(\Gamma\subset\Qbar_{>0}^{\times}\) have multiplicative rank \(r\).
For every finite \(A\subset\Gamma\),
\begin{equation}\label{eq:E1-finite-rank-defect}
 \boxed{\displaystyle
 \trdeg_KK(E_1(\alpha):\alpha\in A)\ge |A|-r-1.}
\end{equation}
There is a set \(E\subset\Gamma\), \(|E|\le r+1\), such that
\begin{equation}\label{eq:E1-cofinite}
 \{E_1(\alpha):\alpha\in\Gamma\setminus E\}
\end{equation}
is algebraically independent over \(K\).  In particular at most \(r+1\)
members of the full tail family are algebraic over \(K\), and at most
\(r+1\) members of
\[
 \{\gamma\}\cup\{E_1(\alpha):\alpha\in\Gamma\}
\]
are algebraic over \(K\).  If \(\gamma\in K^{\mathrm{alg}}\), every
occurrence of \(r+1\) in these tail assertions can be replaced by \(r\).
\end{theorem}

\begin{proof}
By \cref{thm:cofinite-independence},
\[
 |A|-r\le\trdeg_KK(c_\alpha:\alpha\in A).
\]
Since \(c_\alpha=\gamma+E_1(\alpha)\), the field on the right is contained
in \(K(\gamma,E_1(\alpha):\alpha\in A)\).  Adjoining \(\gamma\) costs at
most one transcendence degree, which proves
\eqref{eq:E1-finite-rank-defect}.  Choose a maximal algebraically
independent subfamily of the tails.  If its complement contained \(r+2\)
elements, adjoining the finitely many members of the maximal family needed
to algebraize them would contradict \eqref{eq:E1-finite-rank-defect}.
This proves \eqref{eq:E1-cofinite} and the first algebraicity count.

If \(\gamma\) is algebraic over \(K\), adjoining it costs nothing and the
same proof gives defect \(r\).  If it is transcendental, it is not counted
among the algebraic members.  The assertion for the union follows in both
cases.
\end{proof}

\begin{corollary}[The two-prime triangular grid]
\label{cor:two-prime-grid}
For \(n\ge1\), put
\[
 A_n=\{2^k3^\ell:k,\ell\ge0,\ 1\le k+\ell\le n\},
 \qquad |A_n|=\frac{n(n+3)}2.
\]
Among \(\{c_\alpha:\alpha\in A_n\}\), all but at most two can be chosen
jointly algebraically independent over \(K\), and at most two are
algebraic.  Among \(\{E_1(\alpha):\alpha\in A_n\}\), all but at most three
can be chosen jointly algebraically independent, and at most three are
algebraic.  Thus at \(n=20\) the grid has \(230\) entries: at least \(228\)
of the first family, and at least \(227\) of the tail family, belong to
jointly algebraically independent subfamilies.
\end{corollary}

\begin{proof}
The group generated by \(2\) and \(3\) has multiplicative rank two.  Apply
\cref{thm:cofinite-independence,thm:E1-cofinite}; the cardinality is
\(\sum_{j=1}^n(j+1)=n(n+3)/2\).
\end{proof}

\begin{corollary}[Four classical special-function families]
\label{cor:four-classical-cofinite}
Let \(\Gamma\subset\Qbar_{>0}^{\times}\) have multiplicative rank \(r\).
From the combined family
\begin{equation}\label{eq:four-classical-family}
 \{E_1(\alpha),\Ei(\alpha),\Ci(\alpha),
   \operatorname{Si}(\alpha):\alpha\in\Gamma\}
\end{equation}
delete at most \(r+1\) members.  The remaining family is algebraically
independent over \(K\); in particular at most \(r+1\) members of the full
family are algebraic over \(K\).

For \(\alpha=1,\ldots,M\), the \(4M\) displayed values have transcendence
degree at least
\begin{equation}\label{eq:four-classical-integer-count}
 4M-\pi(M)-1
\end{equation}
over \(K\), and at most \(\pi(M)+1\) of them are algebraic over \(K\).
\end{corollary}

\begin{proof}
For each positive \(\alpha\), the four independent coordinates
\[
 \Ein(\alpha),\quad \Ein(-\alpha),\quad
 \frac{\Ein(i\alpha)+\Ein(-i\alpha)}2,\quad
 \frac{\Ein(i\alpha)-\Ein(-i\alpha)}{2i}
\]
are respectively
\[
 \gamma+\log\alpha+E_1(\alpha),\quad
 \gamma+\log\alpha-\Ei(\alpha),\quad
 \gamma+\log\alpha-\Ci(\alpha),\quad
 \operatorname{Si}(\alpha).
\]
Across distinct positive \(\alpha\), the underlying points
\(\{\pm\alpha,\pm i\alpha\}\) are disjoint; hence all these coordinates
are algebraically independent over \(K\) by
\cref{thm:input-independence,lem:linear-images}.  For any finite selection
from \eqref{eq:four-classical-family}, adjoining \(\gamma\) and at most
\(r\) logarithmic generators recovers the same number of independent
coordinates.  The resulting defect is at most \(r+1\), and maximality gives
the cofinite claim.  The integers up to \(M\) generate a multiplicative
group of rank \(\pi(M)\), which proves
\eqref{eq:four-classical-integer-count}.
\end{proof}

The finite-dimensional form identifies exactly where the deficit occurs.
Let \(\alpha_1,\ldots,\alpha_n\) be positive algebraic numbers contained in
a rank-\(r\) group.  Choose logarithmic generators
\(\lambda=(\log g_1,\ldots,\log g_r)^{\mathsf T}\) and an exponent matrix
\(L\in M_{n\times r}(\Q)\) such that
\begin{equation}\label{eq:log-exponent-matrix}
 (\log\alpha_1,\ldots,\log\alpha_n)^{\mathsf T}=L\lambda.
\end{equation}

\begin{theorem}[Linear-image rank deficit]
\label{thm:linear-rank-deficit}
Let \(M\in M_{m\times n}(\Qbar)\) have row rank \(s\), and put
\(t=\operatorname{rank}(ML)\).  With
\(c=(c_{\alpha_1},\ldots,c_{\alpha_n})^{\mathsf T}\), one has
\begin{equation}\label{eq:linear-rank-deficit}
 \boxed{\displaystyle
 \trdeg_KK(Mc)\ge s-t\ge s-r.}
\end{equation}
\end{theorem}

\begin{proof}
Write \(x=(\Ein(\alpha_1),\ldots,\Ein(\alpha_n))^{\mathsf T}\).  Then
\(x=c+L\lambda\).  Choose a row selector
\(R\in M_{s\times m}(\Qbar)\) for which \(RM\) has row rank \(s\).
The \(s\) coordinates of \(RMx\) are algebraically independent by
\cref{lem:linear-images}.  On the other hand,
\[
 RMx=R(Mc)+(RML)\lambda,
 \qquad \operatorname{rank}(RML)\le\operatorname{rank}(ML)=t.
\]
Thus \(\trdeg_KK((RML)\lambda)\le t\), while \(R(Mc)\) is a linear image
of \(Mc\).  Therefore
\[
 s\le \trdeg_KK(Mc)+t,
\]
which is the assertion.
\end{proof}

There is also a sharp statement in the directions where both the Euler
coefficient and all logarithms cancel.

\begin{corollary}[Complete cancellation space]
\label{cor:complete-cancellation-space}
Let
\begin{equation}\label{eq:cancellation-space}
 U=\ker\begin{pmatrix}\mathbf1^{\mathsf T}\\L^{\mathsf T}\end{pmatrix}
 \subset\Qbar^n,
 \qquad
 d=\dim U=n-\operatorname{rank}[\mathbf1,L].
\end{equation}
For any basis \(u_1,\ldots,u_d\) of \(U\), the numbers
\begin{equation}\label{eq:cancellation-values}
 W_j=\sum_{i=1}^n u_{j,i}\Ein(\alpha_i)
     =\sum_{i=1}^n u_{j,i}E_1(\alpha_i)
 \qquad(1\le j\le d)
\end{equation}
are algebraically independent over \(K\).  Thus
\eqref{eq:cancellation-space} is not only the complete space of
gamma-free, log-free coefficient directions: every independent direction
produces a new algebraically independent period.
\end{corollary}

\begin{proof}
The two defining equations for \(U\), together with
\eqref{eq:Ein-connection}, prove the equality in
\eqref{eq:cancellation-values}.  The basis rows have rank \(d\); apply
\cref{lem:linear-images} to the algebraically independent \(\Ein\)-coordinates.
\end{proof}

For a directed graph \(G=(V,E)\), let \(\partial\) be its incidence matrix,
and attach a positive algebraic number \(\alpha_v\) to each vertex.  If
\(L\) is the corresponding exponent matrix, then
\cref{thm:linear-rank-deficit} gives the graph form
\begin{equation}\label{eq:graph-rank-deficit}
 \trdeg_KK\bigl(E_1(\alpha_{t(e)})-E_1(\alpha_{s(e)}):e\in E\bigr)
 \ge |V|-c(G)-\operatorname{rank}(\partial^{\mathsf T}L),
\end{equation}
where \(c(G)\) is the number of connected components.  This is a genuine
incidence-geometric extension of the one-dimensional prime-defect counts
used below.

\section{External realizations of the independence engine}
\label{sec:external-realizations}

We record three consequences that export the algebraic-point input beyond
the exponential-integral language.  They classify relations rather than
asserting new irrationality results for Euler's constant.

\subsection*{Sine-integral graphic matroids}

Let \(G=(V,E)\) be a finite loopless directed multigraph, and choose nonzero algebraic numbers
\(a_v\) such that \(a_u^2\ne a_v^2\) for \(u\ne v\).  Put
\[
 s_v=\operatorname{Si}(a_v)
 =\frac{\Ein(ia_v)-\Ein(-ia_v)}{2i},
 \qquad
 I_e=s_{t(e)}-s_{s(e)}
 =\int_{a_{s(e)}}^{a_{t(e)}}\frac{\sin z}{z}\,dz.
\]

\begin{theorem}[Complete graphic relation ideal]
\label{thm:sine-graphic-ideal}
Let \(\partial_G:K^E\to K^V\) be the incidence map
\(\partial_Ge=e_{t(e)}-e_{s(e)}\).  Then
\begin{equation}\label{eq:sine-graphic-ideal}
 \boxed{\displaystyle
 \ker\!\left(K[X_e:e\in E]\longrightarrow K[I_e:e\in E]\right)
 =
 \left\langle\sum_{e\in E}c_eX_e:
 c\in\ker\partial_G\right\rangle.}
\end{equation}
Consequently
\begin{equation}\label{eq:sine-graphic-trdeg}
 \trdeg_KK(I_e:e\in E)=|V|-c(G),
\end{equation}
and a set of edge integrals is algebraically independent exactly when its
edges form a forest.
\end{theorem}

\begin{proof}
The pairs \(\{ia_v,-ia_v\}\) are disjoint.  The full-rank linear-image
lemma applied to the corresponding algebraically independent \(\Ein\)
values shows that the \(s_v\) are algebraically independent.  The edge
vector is \(I=\partial_G^{\mathsf T}s\).  Hence its coordinate ring is
the symmetric algebra of the incidence row space, whose defining ideal is
generated by the linear cycle relations.  The rank of the incidence matrix
is \(|V|-c(G)\), and a set of incidence columns is independent exactly when
the corresponding edges form a forest in the underlying multigraph.
\end{proof}

This polynomial-relation statement strengthens the disjoint-endpoint
linear-independence results for sine-integral values in
\cite{Delaygue2022} by using the stronger input from the companion paper.

\subsection*{Generalized Dickman log-Laplace values}

For the generalized Dickman distribution of
\cite{GrahovacEtAl2024}, the Laplace transform is
\begin{equation}\label{eq:Dickman-Laplace}
 \psi_{\theta,a}(s)=\exp\{-\theta\Ein(as)\}.
\end{equation}
Take positive algebraic \(\theta_j,a_j,s_j\), set \(x_j=a_js_j\), and
\[
 \Lambda_j=-\log\psi_{\theta_j,a_j}(s_j)=\theta_j\Ein(x_j).
\]

\begin{theorem}[Dickman log-Laplace classification]
\label{thm:Dickman-classification}
For every finite \(J\),
\begin{equation}\label{eq:Dickman-trdeg}
 \boxed{\displaystyle
 \trdeg_KK(\Lambda_j:j\in J)=\#\{x_j:j\in J\}.}
\end{equation}
For each distinct \(x\), choose a representative \(j_x\).  The complete
polynomial relation ideal is
\begin{equation}\label{eq:Dickman-relation-ideal}
 \ker\!\left(K[X_j:j\in J]\to K[\Lambda_j:j\in J]\right)
 =
 \left\langle
 \theta_{j_x}X_j-\theta_jX_{j_x}:
 x_j=x,\ j\ne j_x
 \right\rangle.
\end{equation}
Moreover,
for integers \(m_j\),
\begin{equation}\label{eq:Dickman-multiplicative}
 \boxed{\displaystyle
 \prod_{j\in J}\psi_{\theta_j,a_j}(s_j)^{m_j}=1
 \Longleftrightarrow
 \sum_{j:x_j=x}m_j\theta_j=0
 \quad\text{for every distinct }x.}
\end{equation}
\end{theorem}

\begin{proof}
Choose one representative for each distinct \(x_j\).  The corresponding
\(\Ein(x_j)\) are algebraically independent; every \(\Lambda_j\) is a nonzero
algebraic multiple of its block representative.  This proves the
transcendence degree and the linear defining ideal.  The Laplace values are
positive real numbers, so a multiplicative relation equal to one is
equivalent, after the real logarithm, to a linear relation among the
distinct \(\Ein(x)\); algebraic independence forces the displayed block
conditions.
\end{proof}

The theorem classifies the logarithmic values and the multiplicative
relations equal to one.  It does not claim to classify every polynomial
relation among the Laplace values themselves.

\subsection*{Regularized incomplete-Gamma jets}

For \(\alpha>0\) algebraic, let
\[
 \Gamma_{<}(s,\alpha)=\int_0^\alpha e^{-t}t^{s-1}\,dt,
 \qquad
 \Gamma_{>}(s,\alpha)=\int_\alpha^\infty e^{-t}t^{s-1}\,dt,
\]
initially in their half-planes of convergence and then by continuation.
At \(s=0\), define
\begin{equation}\label{eq:incomplete-gamma-jets}
 R_{\alpha,k}=[s^k]\{s^{-1}-\Gamma_{<}(s,\alpha)\},
 \qquad
 U_{\alpha,k}=[s^k]\Gamma_{>}(s,\alpha)
 \quad(k\ge0).
\end{equation}
This boundary-Laurent and multipoint regime is distinct from the
positive-parameter incomplete-Gamma special values studied, for example,
in \cite{MurtySaha2015}; no priority claim beyond the precise statements
proved below is intended.

\begin{theorem}[Cofinite independence of regularized lower jets]
\label{thm:lower-gamma-jets}
Let \(\Gamma\subset\Qbar_{>0}^{\times}\) have multiplicative rank \(r\),
and let \(S\subset\Gamma\times\N_0\) be finite.  Then
\begin{equation}\label{eq:lower-gamma-jet-defect}
 \boxed{\displaystyle
 \trdeg_KK(R_{\alpha,k}:(\alpha,k)\in S)\ge |S|-r.}
\end{equation}
Consequently, after deleting at most \(r\) members, the entire two-index
family \(\{R_{\alpha,k}:\alpha\in\Gamma,\ k\ge0\}\) is algebraically
independent over \(K\).  In particular at most \(r\) of these values are
algebraic over \(K\).

More explicitly, with \(L_\alpha=\log\alpha\),
\begin{equation}\label{eq:lower-gamma-jet-explicit}
 R_{\alpha,k}=-\frac{L_\alpha^{k+1}}{(k+1)!}
 +\sum_{j=0}^k\frac{(-1)^jL_\alpha^{k-j}}{(k-j)!}
 \Ein_{j+1}(\alpha).
\end{equation}
Thus \(R_{\alpha,0}=c_\alpha\), while
\begin{equation}\label{eq:lower-gamma-at-one}
 R_{1,k}=(-1)^k\Ein_{k+1}(1)
\end{equation}
and the latter all-order family is algebraically independent over \(K\)
without exceptions.
\end{theorem}

\begin{proof}
Splitting \(e^{-t}=1+(e^{-t}-1)\), expanding \(\alpha^s\), and then
expanding \((n+s)^{-1}\) gives \eqref{eq:lower-gamma-jet-explicit}.
For the finite set \(S\), collect the finitely many coordinates
\(\Ein_j(\alpha)\) that occur on the right.  The rows defining the
\(R_{\alpha,k}\) have rank \(|S|\): within each fixed \(\alpha\), the row
of order \(k\) has the new highest coordinate \(\Ein_{k+1}(\alpha)\) with
nonzero coefficient, and different \(\alpha\)-blocks are disjoint.

Let \(X\) be the resulting vector of \(M\) \(\Ein\)-coordinates and let
\(L=K(\lambda_1,\ldots,\lambda_r)\), where the \(\lambda_i\) span the real
logarithms of \(\Gamma\).  Over \(L\), complete the \(|S|\) displayed rows
to an invertible linear coordinate system \((R,Z)\) on \(X\), with
\(|Z|=M-|S|\).  By \eqref{eq:resonant-input},
\[
 M=\trdeg_KK(X)
 \le\trdeg_KL(R,Z)
 \le r+\trdeg_KK(R)+M-|S|.
\]
This proves \eqref{eq:lower-gamma-jet-defect}.  The maximality argument of
\cref{thm:cofinite-independence} then gives the cofinite assertion.
Setting \(L_1=0\) proves \eqref{eq:lower-gamma-at-one}.
\end{proof}

Write \(\Gamma(s)=s^{-1}+\sum_{k\ge0}g_ks^k\).  Since
\(\Gamma=\Gamma_<+\Gamma_>\), one has
\begin{equation}\label{eq:upper-lower-jet-split}
 U_{\alpha,k}=g_k+R_{\alpha,k}.
\end{equation}

\begin{corollary}[Upper-jet defect]
\label{cor:upper-gamma-jets}
For finite \(S\subset\Gamma\times\N_0\), let
\(d(S)=\#\{k:(\alpha,k)\in S\}\).  Then
\begin{equation}\label{eq:upper-gamma-jet-defect}
 \boxed{\displaystyle
 \trdeg_KK(U_{\alpha,k}:(\alpha,k)\in S)
 \ge |S|-r-d(S).}
\end{equation}
For every fixed \(k\), all but at most \(r+1\) members of
\(\{U_{\alpha,k}:\alpha\in\Gamma\}\) form an algebraically independent
family over \(K\).  At \(k=0\), \(U_{\alpha,0}=E_1(\alpha)\), so this
recovers \cref{thm:E1-cofinite}.
\end{corollary}

\begin{proof}
By \eqref{eq:upper-lower-jet-split}, adjoining the \(d(S)\) Gamma
coefficients \(g_k\) to the upper jets generates all lower jets in \(S\).
Combine \eqref{eq:lower-gamma-jet-defect} with the transcendence-degree
inequality.  For fixed \(k\), apply the same maximality argument as before.
\end{proof}

The finite-order shift ledger can now be replaced by a complete invariant
ring.  This is a structural classification, not an irrationality theorem.

\begin{proposition}[Complete finite Gamma-jet quotient]
\label{prop:finite-gamma-jet-quotient}
Fix \(N\ge0\), let
\[
 A_N(s)=1+\sum_{k=0}^Ng_ks^{k+1}\pmod {s^{N+2}},
 \qquad T_k=U_{1,k},\qquad d_k=g_k-T_k,
\]
and put \(\kappa_j=[s^j]\log A_N(s)\).  On the polynomial ring in the
\(g_k,T_k\), define a \(\mathbb G_a\)-action by
\begin{equation}\label{eq:gamma-jet-shift}
 A_N(s)\longmapsto e^{-ts}A_N(s)\pmod {s^{N+2}},
 \qquad d_k\longmapsto d_k.
\end{equation}
Then
\begin{equation}\label{eq:gamma-jet-invariant-ring}
 \boxed{\displaystyle
 \Q[g_0,\ldots,g_N,T_0,\ldots,T_N]^{\mathbb G_a}
 =\Q[\kappa_2,\ldots,\kappa_{N+1},d_0,\ldots,d_N].}
\end{equation}
The analogous equality holds for fraction fields.  At the actual Gamma
split,
\begin{equation}\label{eq:actual-gamma-jet-invariants}
 \kappa_1=-\gamma,qquad
 \kappa_j=\frac{(-1)^j}{j}\zeta(j)\ (j\ge2),qquad
 d_k=(-1)^{k+1}\Ein_{k+1}(1).
\end{equation}
Thus \(\gamma\) is the unique orbit coordinate at every finite order; the
zeta cumulants and resonant \(\Ein\)-coordinates generate all polynomial
and rational invariants.
\end{proposition}

\begin{proof}
The change of variables
\[
 (g_0,\ldots,g_N,T_0,\ldots,T_N)
 \longleftrightarrow
 (\kappa_1,\ldots,\kappa_{N+1},d_0,\ldots,d_N)
\]
is triangular and polynomial in both directions, by the truncated formal
identities \(\log A_N\) and \(A_N=\exp(\log A_N)\).  Under
\eqref{eq:gamma-jet-shift}, \(\kappa_1\mapsto\kappa_1-t\) and every other
displayed coordinate is fixed.  The invariant ring of translation in one
variable is generated by the remaining variables, proving
\eqref{eq:gamma-jet-invariant-ring}.  Finally
\(A_N(s)\) is the truncation of
\(s\Gamma(s)=\Gamma(1+s)\), so
\[
 \log\Gamma(1+s)=-\gamma s+
 \sum_{j\ge2}\frac{(-1)^j}{j}\zeta(j)s^j.
\]
Equation \eqref{eq:upper-lower-jet-split} at \(\alpha=1\), together with
\eqref{eq:lower-gamma-at-one}, gives the formula for \(d_k\).
\end{proof}

\begin{remark}[Correct reading of the all-order ledger]
The gauge orbit has dimension exactly one at every order.  Any additional
unknown excess in a truncated equation count is not another copy of the
Euler shift: the new odd-zeta cumulants are themselves invariant arithmetic
coordinates.  Consequently a growing unknown-minus-equation count is useful
bookkeeping, but is not by itself a no-go theorem.
\end{remark}

The corresponding export to arbitrary characteristic-zero algebraic
matroids uses an external linearization theorem and is therefore separated
from these direct applications; see \cref{app:Ein-matroids}.

\begin{remark}[A conditional density consequence of Powers]
\label{rem:Powers-conditional}
In the notation of Powers' later multisection manuscript
\cite[\S5]{Powers2025}, assume that
\[
 G_1(1),H_1(1),\ldots,G_N(1),H_N(1)
\]
are algebraically independent over \(\Qbar(e)\) for every \(N\).  This
assertion appears in the concluding multisection discussion rather than as
a proved numbered theorem in the cited version.  There
\begin{align*}
 G_n(1)&=\frac{\widetilde\delta^{(n)}-\delta^{(n)}}{2e},\\
 H_n(1)&=\gamma^{(n)}+
 \frac{\widetilde\delta^{(n)}+\delta^{(n)}}{2e}.
\end{align*}
Conditional on it, combine the two generalized Euler--Gompertz rows
\[
 \delta^{(1)},\ldots,\delta^{(N)},\qquad
 \widetilde\delta^{(1)},\ldots,\widetilde\delta^{(N)}.
\]
Their field satisfies
\[
 \trdeg_{\Qbar}\Qbar\bigl(
 \delta^{(1)},\ldots,\delta^{(N)},
 \widetilde\delta^{(1)},\ldots,\widetilde\delta^{(N)}\bigr)
 \ge\left\lfloor\frac{3N}{2}\right\rfloor-1.
\]
Indeed the Gamma-moment identities quoted by Powers give
\[
 \trdeg_{\Qbar}\Qbar(\gamma^{(1)},\ldots,\gamma^{(N)})
 \le N-\left\lfloor\frac N2\right\rfloor+1.
\]
The assumed \(2N\) independent coordinates, together with the
transcendence of \(e\), lie in the field generated by \(e\), the two
Euler--Gompertz rows, and these Gamma moments.  Subtracting transcendence
degrees yields the displayed bound.  Greedy basis selection in the
associated algebraic matroid, ordered as
\(\delta^{(1)},\widetilde\delta^{(1)},\delta^{(2)},\ldots\), then gives a
jointly algebraically independent subfamily of lower density at least
\(3/4\).  This entire paragraph is conditional, and no unconditional
result elsewhere in the paper depends on it.
\end{remark}

\section{Universal level-trace polynomials}
\label{sec:trace-polynomials}

For \(c\in\C^\times\), set
\[
 F_c(z)=\Ein(z)-c,
\]
and denote by \(Z_c\) its zero multiset, with multiplicities.  Since \(F_c\)
has order one, its inverse-power sums converge absolutely from order two
onward.

\begin{theorem}[Universal level-trace polynomial]
\label{thm:universal-trace}
For \(m\ge2\), define
\begin{equation}\label{eq:Pm-def}
 P_m(X):=-m[z^m]\log\bigl(1-X\Ein(z)\bigr).
\end{equation}
Then \(P_m\in\Q[X]\), its constant term is zero, and it is monic of degree
\(m\).  Moreover
\begin{equation}\label{eq:universal-trace}
 \boxed{\displaystyle
 T_m(c):=\sum_{\rho\in Z_c}\rho^{-m}=P_m(c^{-1}).}
\end{equation}
In particular
\begin{align}
 P_2(X)&=X^2-\frac X2,\label{eq:P2}\\
 P_3(X)&=X^3-\frac34X^2+\frac16X,\label{eq:P3}\\
 P_4(X)&=X^4-X^3+\frac{25}{72}X^2-\frac1{24}X.\label{eq:P4}
\end{align}
\end{theorem}

\begin{proof}
A genus-one Hadamard factorization has the form
\[
 F_c(z)=-c\,e^{a_cz}
 \prod_{\rho\in Z_c}
 \left(1-\frac z\rho\right)e^{z/\rho}.
\]
Taking the logarithm at the origin gives, for every \(m\ge2\),
\[
 [z^m]\log\frac{F_c(z)}{F_c(0)}
 =-\frac1m\sum_{\rho\in Z_c}\rho^{-m}.
\]
On the other hand,
\[
 \frac{F_c(z)}{F_c(0)}=1-\frac{\Ein(z)}c.
\]
Comparison proves \eqref{eq:universal-trace}.  Formula
\eqref{eq:Pm-def} is a polynomial in \(X\) of degree at most \(m\).  Since
\(\Ein(z)=z+O(z^2)\), the contribution of
\(X^m\Ein(z)^m/m\) to \eqref{eq:Pm-def} is \(X^m\).  Thus \(P_m\) is monic
of degree \(m\).  The displayed formulas follow by expanding through degree
four.
\end{proof}

\begin{remark}[Multiplicity and genus]
No simplicity assertion is required: all zeros are counted with their
multiplicities.  The first inverse sum is not used.  At order one the
canonical genus-one exponential factor absorbs the non-absolutely convergent
sum, whereas \eqref{eq:universal-trace} is absolute for \(m\ge2\).
\end{remark}

\begin{corollary}[Transversal trace independence]
\label{cor:transversal}
Let \(F\subset\C\) be a field of characteristic zero, and let
\(\{c_j:j\in J\}\subset\C^\times\) be algebraically independent over \(F\).
Choose one integer \(m_j\ge2\) for each \(j\).  Then
\[
 \{T_{m_j}(c_j):j\in J\}
\]
is algebraically independent over \(F\), in the finite-subset sense.
\end{corollary}

\begin{proof}
Fix a finite set \(I\subset J\), and write
\(x_j=c_j^{-1}\) and \(t_j=P_{m_j}(x_j)\).  Since
\[
 P_{m_j}(x_j)-t_j=0,
\]
the field \(F(x_j:j\in I)\) is algebraic over \(F(t_j:j\in I)\).  Hence
\[
 \trdeg_FF(t_j:j\in I)
 =\trdeg_FF(x_j:j\in I)=|I|.
\]
\end{proof}

\begin{theorem}[Exact even-trace recovery and ring closure]
\label{cor:fixed-level-closure}
For every \(c\ne0\),
\begin{equation}\label{eq:even-trace-recovery}
 \boxed{\displaystyle
 c^{-1}=144T_4(c)-144T_2(c)^2-14T_2(c).}
\end{equation}
Consequently, for every characteristic-zero field \(F\subset\C\),
\[
 \boxed{\displaystyle
 F[T_m(c):m\ge2]
 =F[T_2(c),T_4(c)]
 =F[c^{-1}]}
\]
and
\[
 \boxed{F(T_2(c),T_4(c))=F(c).}
\]
\end{theorem}

\begin{proof}
With \(x=c^{-1}\), \eqref{eq:P2} and \eqref{eq:P4} give
\[
 144P_4(x)-144P_2(x)^2-14P_2(x)=x.
\]
This proves \eqref{eq:even-trace-recovery}.  Every trace is a polynomial in
\(x\) by \cref{thm:universal-trace}, while the displayed identity recovers
\(x\) from \(T_2,T_4\).  Taking fraction fields and using \(F(x)=F(c)\)
finishes the proof.
\end{proof}

The preceding recovery identity also gives the exact, rather than merely
heuristic, form of the finite-stage shift obstruction.  It is important to
distinguish this gauge action, which fixes every connection coordinate,
from common translation of the connection coordinates themselves;
augmentation zero pertains to the latter action, not to the former.

\begin{theorem}[Exact finite-stage shift quotient]
\label{thm:exact-shift-quotient}
Let \(B\) be a characteristic-zero field, put
\[
 R=B[G,H_1,\ldots,H_n],\qquad c_i=G+H_i,
\]
and let \(\mathbb G_a\) act by
\begin{equation}\label{eq:gauge-action}
 G\longmapsto G+t,\qquad H_i\longmapsto H_i-t.
\end{equation}
Then
\begin{equation}\label{eq:shift-invariant-ring}
 \boxed{R^{\mathbb G_a}=B[c_1,\ldots,c_n],\qquad
 \operatorname{Frac}(R)^{\mathbb G_a}=B(c_1,\ldots,c_n).}
\end{equation}
If the universal traces are formed at these levels, then
\begin{align}
 B\bigl(T_m(c_i):1\le i\le n,\ m\ge2\bigr)
 &=B(c_1,\ldots,c_n),\label{eq:trace-invariant-field}\\
 B\bigl[T_m(c_i):1\le i\le n,\ m\ge2\bigr]
 &=B[c_1^{-1},\ldots,c_n^{-1}].\label{eq:trace-invariant-ring}
\end{align}
Thus the complete finite multi-level trace field is exactly the rational
quotient by the gauge direction.  It loses one orbit coordinate and no
others.
\end{theorem}

\begin{proof}
The polynomial change of variables \(H_i=c_i-G\) identifies
\(R=B[c_1,\ldots,c_n,G]\), and \eqref{eq:gauge-action} becomes translation
of \(G\) with every \(c_i\) fixed.  In characteristic zero the polynomial
and rational invariants of translation in one variable are respectively
\(B[c_1,\ldots,c_n]\) and \(B(c_1,\ldots,c_n)\), proving
\eqref{eq:shift-invariant-ring}.  Every trace is a polynomial in
\(c_i^{-1}\), while \eqref{eq:even-trace-recovery} recovers \(c_i^{-1}\)
from \(T_2(c_i),T_4(c_i)\).  This proves
\eqref{eq:trace-invariant-field}--\eqref{eq:trace-invariant-ring}.
\end{proof}

\begin{remark}[What the formal quotient does and does not prove]
The elimination ideal of the equations \(c_i=G+H_i\) in \(B[G]\) is zero:
the same finite invariant data admit a model with \(G=0\) and another with
\(G\) transcendental.  Hence invariant algebraic identities alone cannot
decide the arithmetic of the orbit coordinate.  This is a statement about
the formal axiom system, not about the actual complex number \(\gamma\).
Indeed the analytic completion supplies extra boundary data.  If
\(y_N\to g\) in a Hausdorff topological field and a continuous action fixes
every \(y_N\), it fixes \(g\).  Applied to \(Y_N\to\gamma\), and to the two
carrier limits \(\tau_N\to S_2\), \(\sigma_N\to S_4\), this is precisely
where the finite-stage gauge symmetry breaks; \(S_2,S_4\) recover
\(\gamma\) by \eqref{eq:gamma-even-recovery}.
\end{remark}

\begin{corollary}[Arithmetic classification of every Euler trace]
\label{cor:Euler-trace-arithmetic}
For \(m\ge2\), put \(S_m=T_m(\gamma)\).  Then
\[
 S_m\in\Qbar\quad\Longleftrightarrow\quad\gamma\in\Qbar,
 \qquad
 S_m\text{ is transcendental}
 \quad\Longleftrightarrow\quad\gamma\text{ is transcendental}.
\]
Moreover the two even traces recover Euler's constant without exceptions:
\begin{equation}\label{eq:gamma-even-recovery}
 \boxed{\displaystyle
 \gamma=\frac1{144S_4-144S_2^2-14S_2}.}
\end{equation}
In particular
\[
 \gamma\in\Q\quad\Longleftrightarrow\quad S_2,S_4\in\Q.
\]
\end{corollary}

\begin{proof}
If \(S_m=P_m(\gamma^{-1})\) is algebraic, then \(\gamma^{-1}\) is a root
of the nonconstant polynomial \(P_m(X)-S_m\) over \(\Qbar\), and hence is
algebraic.  The converse is immediate.  Formula
\eqref{eq:gamma-even-recovery} is \eqref{eq:even-trace-recovery} at
\(c=\gamma\); it also proves the rationality equivalence.
\end{proof}

\begin{theorem}[The unique two-trace level collision]
\label{thm:T2T3-collision}
For \(c,d\in\C^\times\),
\[
 (T_2(c),T_3(c))=(T_2(d),T_3(d))
\]
holds if and only if either \(c=d\), or
\begin{equation}\label{eq:T2T3-exceptional-pair}
 \boxed{\{c,d\}=\{6-2\sqrt3,6+2\sqrt3\}.}
\end{equation}
In the exceptional case the common pair is
\begin{equation}\label{eq:T2T3-exceptional-value}
 \boxed{(T_2,T_3)=\left(-\frac1{24},\frac1{96}\right).}
\end{equation}
\end{theorem}

\begin{proof}
Put \(x=c^{-1}\) and \(y=d^{-1}\).  From \eqref{eq:P2}, equality of the
second traces gives
\[
 (x-y)\left(x+y-\frac12\right)=0.
\]
If \(x=y\), then \(c=d\).  Otherwise \(x+y=1/2\), and equality of the
third traces, using \eqref{eq:P3}, reduces to
\[
 x^2+xy+y^2-\frac34(x+y)+\frac16=0,
\]
or \(xy=1/24\).  Thus \(x,y\) are the roots of
\[
 X^2-\frac12X+\frac1{24},
\]
and taking reciprocals gives \eqref{eq:T2T3-exceptional-pair}.  Direct
substitution in \eqref{eq:P2} and \eqref{eq:P3} gives
\eqref{eq:T2T3-exceptional-value}.
\end{proof}

\begin{corollary}[Exact vanishing traces with transcendental divisors]
\label{cor:vanishing-trace-levels}
For every \(m\ge2\), define
\[
 \mathcal V_m=
 \{c\in\C^\times:T_m(c)=0\}.
\]
Then
\begin{equation}\label{eq:vanishing-trace-levels}
 \boxed{\mathcal V_m=
 \{x^{-1}:x\ne0,\ P_m(x)=0\}\subset\Qbar^\times.}
\end{equation}
This set is nonempty and finite; its elements arise from exactly \(m-1\)
nonzero roots of \(P_m\), counted with multiplicity.  For every
\(c\in\mathcal V_m\), every zero \(\rho\) of \(\Ein(z)-c\) is
transcendental, although the absolutely convergent inverse trace satisfies
\[
 \sum_{\Ein(\rho)=c}\rho^{-m}=0.
\]
For example, \(\mathcal V_2=\{2\}\).
\end{corollary}

\begin{proof}
The trace identity \eqref{eq:universal-trace} gives
\eqref{eq:vanishing-trace-levels}.  The polynomial \(P_m\) has degree
\(m\), constant term zero, and linear coefficient
\[
 [X]P_m(X)=m[z^m]\Ein(z)=\frac{(-1)^{m-1}}{m!}\ne0.
\]
Hence zero is a simple root and the remaining \(m-1\) roots are nonzero,
counted with multiplicity.  They are algebraic because \(P_m\in\Q[X]\).

Now fix \(c\in\mathcal V_m\).  If a zero \(\rho\) of \(\Ein(z)-c\) were
algebraic, then \(\rho\ne0\) and
\(\Ein(\rho)=c\in\Qbar\), contradicting
\cref{thm:input-independence}.  Thus every point of the divisor is
transcendental.  Finally \(P_2(X)=X(X-1/2)\) gives the example.
\end{proof}

Thus even all inverse traces at one fixed level still contain only one
coordinate.  Within the universal inverse-trace construction, varying the
trace order at a fixed level produces only the one-dimensional field
generated by that level.  A transversal algebraically independent
inverse-trace family therefore requires varying the level; other constructions
may instead vary a coefficient, a power orbit, or an external parameter.

\section{Balanced dyadic levels and all-order moving traces}
\label{sec:dyadic}

For \(N\ge1\), define
\begin{equation}\label{eq:YN}
 Y_N=(N+1)\Ein(2^N)-N\Ein(2^{N+1}).
\end{equation}
The companion paper proved the first assertion below.  The remaining
assertions are the analytic bridge to the Euler level.

\begin{proposition}[Balanced dyadic approach]
\label{prop:dyadic-approach}
The family \(\{Y_N:N\ge1\}\) is countably algebraically independent over
\(K\).  Moreover
\begin{equation}\label{eq:epsilonN}
 Y_N=\gamma+\varepsilon_N,
 \qquad
 \varepsilon_N=(N+1)E_1(2^N)-NE_1(2^{N+1}),
\end{equation}
and
\begin{equation}\label{eq:epsilon-properties}
 \varepsilon_N>0,\qquad
 \varepsilon_N\downarrow0,\qquad
 \varepsilon_N\sim\frac{N+1}{2^N}e^{-2^N}.
\end{equation}
In particular \(Y_N\downarrow\gamma\).
\end{proposition}

\begin{proof}
The algebraic independence is \cite[Lem.~11.2]{KEIO2026I}; it also follows directly from
\cref{thm:input-independence,lem:linear-images}, because the finite row family
\(\bigl((N+1)e_N-Ne_{N+1}\bigr)\) has full row rank.

Substitution of \eqref{eq:Ein-connection} in \eqref{eq:YN} cancels both the
logarithmic and the \(N\log2\) terms and gives \eqref{eq:epsilonN}.  Positivity
follows from
\[
 \varepsilon_N
 =E_1(2^N)+N\{E_1(2^N)-E_1(2^{N+1})\}>0.
\]
For strict decrease put \(g(t)=E_1(2^t)\).  Since
\[
 g'(t)=-(\log2)e^{-2^t},\qquad
 g''(t)=(\log2)^2 2^t e^{-2^t}>0,
\]
strict convexity gives
\[
 \varepsilon_N-\varepsilon_{N+1}
 =(N+1)\{g(N)-2g(N+1)+g(N+2)\}>0.
\]
Finally, the standard positive-axis expansion
\begin{equation}\label{eq:E1-asymptotic}
 E_1(x)=\frac{e^{-x}}x\left(1+O(x^{-1})\right)
 \qquad(x\to+\infty)
\end{equation}
gives the asymptotic in \eqref{eq:epsilon-properties}; see
\cite[\S6.12\(i\)]{DLMF}.
\end{proof}

For \(m\ge2\), let
\begin{equation}\label{eq:tau-mN}
 \tau_{m,N}:=\sum_{\Ein(\rho)=Y_N}\rho^{-m},
\end{equation}
where the solutions are counted with multiplicity.

\begin{theorem}[All-order transversal tower]
\label{thm:all-order-tower}
For every \(m\ge2\) and \(N\ge1\),
\begin{equation}\label{eq:tau-Pm}
 \tau_{m,N}=P_m(Y_N^{-1}).
\end{equation}
For any selection \(m_N\ge2\), the family
\[
 \{\tau_{m_N,N}:N\ge1\}
\]
is countably algebraically independent over \(K\).  For each fixed
\(m\ge2\),
\begin{equation}\label{eq:tau-limit}
 \tau_{m,N}\longrightarrow S_m:=P_m(\gamma^{-1}).
\end{equation}
\end{theorem}

\begin{proof}
Equation \eqref{eq:tau-Pm} is \cref{thm:universal-trace}.  Apply
\cref{cor:transversal} to the independent levels \(Y_N\).  The limit follows
from \cref{prop:dyadic-approach} and continuity of \(P_m(1/x)\) at
\(x=\gamma\).
\end{proof}

\begin{remark}[Why transversal is sharp]
For a fixed \(N\), the even traces already recover the level exactly:
\begin{equation}\label{eq:YN-even-recovery}
 \boxed{\displaystyle
 Y_N^{-1}=144\tau_{4,N}-144\tau_{2,N}^2-14\tau_{2,N}.}
\end{equation}
Thus
\[
 K[\tau_{m,N}:m\ge2]
 =K[\tau_{2,N},\tau_{4,N}]
 =K[Y_N^{-1}].
\]
The full double family is therefore not algebraically independent: the
theorem gives exactly one independent spectral coordinate per moving level.
\end{remark}

The order-two row has additional positivity.  Put
\begin{equation}\label{eq:f-def}
 f(x)=\frac1{x^2}-\frac1{2x}=\frac{2-x}{2x^2},
 \qquad \tau_N:=\tau_{2,N}=f(Y_N).
\end{equation}

\begin{theorem}[Exact positive convergence at order two]
\label{thm:tau2-convergence}
One has
\[
 0<\tau_1<\tau_2<\cdots<S_2,
 \qquad \tau_N\longrightarrow S_2.
\]
More precisely,
\begin{equation}\label{eq:S2-tail-exact}
 \boxed{\displaystyle
 S_2-\tau_N=
 \frac{\varepsilon_N
 \{\gamma(4-\gamma)+(2-\gamma)\varepsilon_N\}}
 {2\gamma^2(\gamma+\varepsilon_N)^2}.}
\end{equation}
Consequently
\begin{equation}\label{eq:S2-tail-asymptotic}
 S_2-\tau_N\sim
 \frac{4-\gamma}{2\gamma^3}
 \frac{N+1}{2^N}e^{-2^N}.
\end{equation}
\end{theorem}

\begin{proof}
Since \(E_1(2)<e^{-2}/2\) and \(\gamma<1\),
\(Y_1=\gamma+2E_1(2)-E_1(4)<\gamma+e^{-2}<2\), so all \(Y_N\) lie in
\((0,2)\).  But
\[
 f'(x)=\frac{x-4}{2x^3}<0\qquad(0<x<2).
\]
Thus \(Y_N\downarrow\gamma\) implies \(\tau_N\uparrow f(\gamma)=S_2\), and
positivity follows from \(f(x)>0\) on \((0,2)\).  Direct subtraction with
\(Y_N=\gamma+\varepsilon_N\) gives \eqref{eq:S2-tail-exact}.  Combine it
with \eqref{eq:epsilon-properties} to obtain
\eqref{eq:S2-tail-asymptotic}.
\end{proof}

\begin{theorem}[Non-P-recursiveness of the dyadic moving sequences]
\label{thm:non-p-recursive}
Neither \((Y_N)_{N\ge1}\) nor \((\tau_N)_{N\ge1}\) is P-recursive over
\(\C\).  Equivalently, neither sequence satisfies a nonzero recurrence
\[
 \sum_{j=0}^r p_j(N)u_{N+j}=0
 \qquad\bigl(p_j\in\C[N]\bigr)
\]
for all sufficiently large \(N\).
\end{theorem}

\begin{proof}
Suppose first that such a recurrence holds for \(Y_N=\gamma+\varepsilon_N\).
Since every polynomial multiple of \(\varepsilon_{N+j}\) tends to zero, we
obtain
\[
 \gamma\sum_{j=0}^rp_j(N)\longrightarrow0.
\]
The sum is a polynomial, hence it vanishes identically.  Let \(j_0\) be the
least index for which \(p_{j_0}\ne0\).  Dividing the remaining relation by
\(p_{j_0}(N)\varepsilon_{N+j_0}\) gives
\[
 1+\sum_{j>j_0}
 \frac{p_j(N)}{p_{j_0}(N)}
 \frac{\varepsilon_{N+j}}{\varepsilon_{N+j_0}}=0.
\]
But \eqref{eq:epsilon-properties} implies, for every fixed \(j>j_0\),
\[
 \frac{\varepsilon_{N+j}}{\varepsilon_{N+j_0}}
 =\exp\{-\bigl(2^j-2^{j_0}\bigr)2^N+O(N)\},
\]
which dominates every rational function of \(N\).  The displayed equality
therefore tends to \(1=0\), a contradiction.

For \(\tau_N=f(Y_N)\), \cref{thm:tau2-convergence} gives
\[
 \tau_N=S_2+f'(\gamma)\varepsilon_N+O(\varepsilon_N^2),
 \qquad f'(\gamma)\ne0.
\]
The same least-index argument applies verbatim.
\end{proof}

\begin{remark}[Scope of the non-P-recursiveness result]
The theorem does not say that no holonomic construction can involve the
finite-stage values.  It says that the natural moving index itself cannot be
packaged as a fixed P-recursive sequence.  Any application of arithmetic
holonomy must therefore replace the moving tower by a genuinely different,
fixed differential or recurrence system rather than merely reindexing
\(Y_N\).
\end{remark}

For orientation, the first values are shown in \cref{tab:dyadic-values}.
The table is illustrative only; no numerical value is used in a proof.

\begin{table}[ht]
\centering
\caption{The dyadic level and its order-two zero trace.}
\label{tab:dyadic-values}
\begin{tabular}{@{}rcc@{}}
\toprule
\(N\)&\(Y_N\)&\(\tau_N\)\\
\midrule
1&0.671237333907806&1.474569389580200\\
2&0.588478390885392&2.037963940792663\\
3&0.577366307471447&2.133831985200994\\
4&0.577215698103968&2.135171685377284\\
5&0.577215664901535&2.135171980841625\\
\midrule
limit&0.577215664901533&2.135171980841646\\
\bottomrule
\end{tabular}
\end{table}

\section{Positive spectral carriers for
\texorpdfstring{\(S_2\) and \(S_4\)}{S2 and S4}}
\label{sec:spectral}

Define
\begin{equation}\label{eq:dN-def}
 d_0=\tau_1,
 \qquad
 d_N=\tau_{N+1}-\tau_N\quad(N\ge1).
\end{equation}

\begin{theorem}[Positive algebraically independent increments]
\label{thm:increments}
The sequence \(\{d_N:N\ge0\}\) is positive and countably algebraically
independent over \(K\).  It satisfies
\begin{equation}\label{eq:S2-sum}
 \boxed{\displaystyle S_2=\sum_{N=0}^\infty d_N}
\end{equation}
and, for \(N\ge1\),
\begin{equation}\label{eq:dN-asymptotic}
 d_N\sim
 \frac{4-\gamma}{2\gamma^3}
 \frac{N+1}{2^N}e^{-2^N}.
\end{equation}
\end{theorem}

\begin{proof}
Positivity follows from \cref{thm:tau2-convergence}.  The finite
transformation
\[
 (\tau_1,\ldots,\tau_{M+1})
 \longleftrightarrow(d_0,\ldots,d_M)
\]
is an invertible triangular \(\Q\)-linear transformation.  The
\(\tau_N\) are algebraically independent by
\cref{thm:all-order-tower}; hence so are the \(d_N\).  The telescoping
identity
\[
 \sum_{N=0}^Md_N=\tau_{M+1}
\]
and \cref{thm:tau2-convergence} give \eqref{eq:S2-sum}.  Finally the ratio
of two successive tails in \eqref{eq:S2-tail-asymptotic} tends to zero, so
\[
 d_N=(S_2-\tau_N)-(S_2-\tau_{N+1})\sim S_2-\tau_N,
\]
which proves \eqref{eq:dN-asymptotic}.
\end{proof}

Let \(\{e_N:N\ge0\}\) be the standard basis of \(\ell^2(\N_0)\).

\begin{corollary}[A positive operator carrier]
\label{cor:operator}
The diagonal operator
\[
 \mathsf A e_N=d_Ne_N
\]
is positive, self-adjoint and trace class, with
\[
 \Tr\mathsf A=S_2.
\]
More strongly,
\[
 \sum_{N\ge0}d_N^p<\infty\qquad(p>0).
\]
Thus \(\mathsf A\in\mathcal S_p\) for every \(p\ge1\), and satisfies the
corresponding quasi-Schatten condition for \(0<p<1\).  Its nonzero
eigenvalues are simple and countably algebraically independent over \(K\).
\end{corollary}

\begin{proof}
All assertions except simplicity follow immediately from
\cref{thm:increments}.  The double-exponential asymptotic
\eqref{eq:dN-asymptotic} gives summability of every positive power.  If
\(d_j=d_k\) for \(j\ne k\), the algebraically independent family would
satisfy a nonzero linear relation; hence the eigenvalues are distinct.
\end{proof}

The Fredholm determinant of the operator gives a scalar entire realization.

\begin{theorem}[Order-zero Laguerre--P\'olya realization]
\label{thm:Delta2}
The product
\begin{equation}\label{eq:Delta2}
 \Delta_2(z)=
 \det\!\left(I-\frac{z^2}{2}\mathsf A\right)
 =\prod_{N=0}^\infty\left(1-\frac{d_N}{2}z^2\right)
\end{equation}
converges locally uniformly and defines an order-zero entire function in the
Laguerre--P\'olya class.  Its zeros are simple and equal to
\[
 \pm r_N,\qquad r_N=\sqrt{\frac2{d_N}}.
\]
The positive half \(\{r_N:N\ge0\}\) is countably algebraically independent
over \(K\), and
\begin{equation}\label{eq:Delta2-trace}
 \boxed{\displaystyle
 \sum_{\rho\in Z(\Delta_2)}\rho^{-2}=S_2.}
\end{equation}
\end{theorem}

\begin{proof}
Since \(\sum d_N<\infty\), the product in \eqref{eq:Delta2} converges
normally on compact sets.  Every finite partial product has only real zeros,
so their locally uniform limit belongs to the Laguerre--P\'olya class; see
\cite[Chapter~VIII]{Levin1996}.  For every \(\eta>0\),
\[
 \sum_{N\ge0}r_N^{-\eta}
 =2^{-\eta/2}\sum_{N\ge0}d_N^{\eta/2}<\infty.
\]
Thus the zero exponent of convergence, and hence the order of the canonical
product, is zero.  Simplicity follows from the distinctness of the \(d_N\).

For a finite index set \(I\), the relations \(d_N=2/r_N^2\) imply
\[
 \trdeg_KK(r_N:N\in I)
 \ge \trdeg_KK(d_N:N\in I)=|I|,
\]
which proves algebraic independence of the positive zeros.  Finally,
absolute convergence and \eqref{eq:S2-sum} give
\[
 \sum_{\rho\in Z(\Delta_2)}\rho^{-2}
 =\sum_{N\ge0}\left(r_N^{-2}+(-r_N)^{-2}\right)
 =\sum_{N\ge0}d_N=S_2.
\]
\end{proof}

\begin{theorem}[A second positive carrier and even-trace recovery]
\label{thm:positive-S4-carrier}
Put
\[
 \sigma_N=T_4(Y_N)=P_4(Y_N^{-1})\quad(N\ge1),
\]
and define
\[
 d_0^{(4)}=\sigma_1,\qquad
 d_N^{(4)}=\sigma_{N+1}-\sigma_N\quad(N\ge1).
\]
Then every \(d_N^{(4)}\) is positive, the family
\(\{d_N^{(4)}:N\ge0\}\) is algebraically independent over \(K\), and
\begin{equation}\label{eq:S4-positive-sum}
 S_4=\sum_{N\ge0}d_N^{(4)}.
\end{equation}
Consequently the diagonal operator
\(\mathsf A_4=\diag(d_N^{(4)}:N\ge0)\) belongs to every Schatten class and
\(\Tr\mathsf A_4=S_4\).  Together with the order-two carrier
\(\mathsf A_2:=\mathsf A\) of \cref{cor:operator}, it gives
\begin{equation}\label{eq:gamma-two-positive-carriers}
 \boxed{\displaystyle
 \gamma=
 \left\{144\Tr\mathsf A_4
 -144(\Tr\mathsf A_2)^2-14\Tr\mathsf A_2\right\}^{-1}.}
\end{equation}
\end{theorem}

\begin{proof}
The earlier bound \(Y_N<2\) and the monotonicity \(Y_N\downarrow\gamma\)
show that \(x_N=Y_N^{-1}\) increases and satisfies \(x_N>1/2\).  From
\eqref{eq:P4},
\[
 P_4(1/2)=\frac1{288},\qquad P_4'(1/2)=\frac1{18},
\]
and
\[
 P_4''(x)=12x^2-6x+\frac{25}{36}>0\qquad(x\ge1/2).
\]
Thus \(P_4\) is positive and strictly increasing on this range, so all
increments are positive.  The \(\sigma_N\) are algebraically
independent by \cref{thm:all-order-tower}; the triangular change from
\((\sigma_1,\ldots,\sigma_{M+1})\) to
\((d_0^{(4)},\ldots,d_M^{(4)})\) preserves this property.  Telescoping and
\(T_4(Y_N)\to T_4(\gamma)=S_4\) prove
\eqref{eq:S4-positive-sum}.  The doubly exponential convergence of \(Y_N\)
gives summability of every positive power of \(d_N^{(4)}\).  Finally apply
\eqref{eq:gamma-even-recovery}.
\end{proof}

\begin{remark}[Three distinct meanings of ``gamma-free'']
Each finite datum \(Y_N,\tau_N,d_N\) is defined by algebraic-point
\(\Ein\)-values without displaying \(\gamma\).  The limit and sum
\[
 S_2=\lim_N\tau_N=\sum_Nd_N
\]
are countable constructions.  They do not prove that \(S_2\) belongs to the
finite field generated by algebraic-point \(\Ein\)-values.  Indeed
\([z^2]\Delta_2(z)=-S_2/2\), so its Taylor coefficients are already beyond
what the finite-stage algebraic-independence theorem controls.  No claim is
made that \(\Delta_2\) is an \(E\)-function, a \(D\)-finite function, or an
arithmetic entire function.
\end{remark}

\begin{remark}[No arithmetic conclusion from the infinite sum]
Algebraic independence controls finite polynomial relations.  It says
nothing by itself about the arithmetic nature of an infinite sum.  Thus
\eqref{eq:S2-sum} and \eqref{eq:Delta2-trace} do not prove that \(S_2\), or
equivalently \(\gamma\), is transcendental.
\end{remark}

\section{Further moving-level towers}
\label{sec:other-towers}

The dyadic tower is not isolated.  This section gives a fixed-coefficient
power-orbit family and a logarithmic lattice.  They have different decay
rates and different independence defects.

\subsection{Fixed-coefficient power orbits}

Fix an integer \(n\ge2\) and a positive algebraic number \(q>1\).  For
\(k\ge0\), put
\begin{equation}\label{eq:Bnk}
 B_{n,k}(q)=
 \frac{n\Ein(q^k)-\Ein(q^{nk})}{n-1}.
\end{equation}

\begin{theorem}[Power-orbit moving tower]
\label{thm:power-orbit}
For fixed \(n\) and \(q\), the family
\(\{B_{n,k}(q):k\ge0\}\) is countably algebraically independent over \(K\).
Moreover
\begin{equation}\label{eq:eta-nk}
 B_{n,k}(q)=\gamma+\eta_{n,k}(q),\qquad
 \eta_{n,k}(q)=
 \frac{nE_1(q^k)-E_1(q^{nk})}{n-1},
\end{equation}
and
\begin{equation}\label{eq:eta-properties}
 \eta_{n,k}>0,\qquad
 \eta_{n,k}\downarrow0,\qquad
 \eta_{n,k}\sim
 \frac n{n-1}\frac{e^{-q^k}}{q^k}.
\end{equation}
For every choice \(m_k\ge2\), the exact traces
\[
 T_{m_k}(B_{n,k}(q))
\]
are countably algebraically independent over \(K\).  For fixed \(m\), they
converge to \(S_m\).  At \(m=2\) they increase to \(S_2\) and yield, by
successive differences, another positive algebraically independent spectral
realization of \(S_2\).
\end{theorem}

\begin{proof}
Let \(X_j=\Ein(q^j)\).  The distinct algebraic points \(q^j\) give an
algebraically independent coordinate family by
\cref{thm:input-independence}.  The row vectors
\[
 v_k=ne_k-e_{nk}\qquad(k\ge0)
\]
are linearly independent: in a finite relation, the coordinate \(nk_0\)
associated with the largest occurring \(k_0\) appears as the terminal
coordinate of that row and as the initial coordinate of no row in the
finite set.  Thus its coefficient vanishes, and descending induction
finishes the proof.  If the finite set consists only of \(k_0=0\), then
\(v_0=(n-1)e_0\ne0\), which is the same conclusion.  Apply
\cref{lem:linear-images}.

The connection formula \eqref{eq:Ein-connection} cancels the logarithmic
terms and gives \eqref{eq:eta-nk}.  For \(x>1\),
\[
 \frac d{dx}\left(
 \frac{nE_1(x)-E_1(x^n)}{n-1}\right)
 =\frac n{n-1}\frac{e^{-x^n}-e^{-x}}x<0.
\]
This proves strict decrease; positivity follows from
\(E_1(x)>E_1(x^n)>0\), and the asymptotic follows from
\eqref{eq:E1-asymptotic}.  The trace assertions follow from
\cref{cor:transversal,thm:universal-trace}.  The order-two construction is
the proof of \cref{thm:increments,thm:Delta2} applied to this decreasing
level sequence: indeed \(B_{n,0}(q)=\Ein(1)<1<2\), since
\((1-e^{-t})/t<1\) on \((0,1]\).
\end{proof}

This tower samples the rank-one witness of the companion paper exactly:
\begin{equation}\label{eq:B-Gn}
 B_{n,k}(q)-\gamma=\frac{G_n(q^k)}{n-1}.
\end{equation}

\begin{remark}[Fixing \(n\) is essential]
The full two-parameter family is not algebraically independent.  For
\(r\ge1\),
\begin{equation}\label{eq:cross-n-relation}
 B_{n^r,k}(q)=\frac{n-1}{n^r-1}
 \sum_{j=0}^{r-1}n^{r-1-j}B_{n,n^jk}(q).
\end{equation}
For example \(B_{4,k}=(2B_{2,k}+B_{2,2k})/3\).
\end{remark}

\subsection{Algebraic-level inverse transcendence}

The first difference between two members of the \(n=2\) power orbit gives
the gamma-free function
\begin{equation}\label{eq:scriptR-def}
 \mathscr R(q)=E_1(1)-2E_1(q)+E_1(q^2)\qquad(q>1).
\end{equation}
It has an unexpectedly rigid inverse on algebraic levels.

\begin{theorem}[Transcendental inverse levels]
\label{thm:inverse-levels}
The map \(\mathscr R:(1,\infty)\to(0,E_1(1))\) is a strictly increasing
bijection.  If \(a\in\Qbar\cap(0,E_1(1))\), then the unique number
\(r_a>1\) satisfying \(\mathscr R(r_a)=a\) is transcendental.

In particular, for every integer \(m\ge5\), the unique root \(r_m>1\) of
\begin{equation}\label{eq:rm-level}
 E_1(1)-2E_1(r_m)+E_1(r_m^2)=\frac1m
\end{equation}
is transcendental, and
\begin{equation}\label{eq:rm-asymptotic}
 r_m=1+\sqrt{\frac em}+\frac e{2m}+O(m^{-3/2}).
\end{equation}
\end{theorem}

\begin{proof}
Direct differentiation gives
\[
 \mathscr R'(q)=\frac2q(e^{-q}-e^{-q^2})>0.
\]
Moreover \(\mathscr R(1)=0\) and
\(\mathscr R(q)\to E_1(1)\) as \(q\to\infty\), proving the bijection.
If \(r_a\) were algebraic, the connection formula would turn
\(\mathscr R(r_a)=a\) into
\[
 \Ein(1)-2\Ein(r_a)+\Ein(r_a^2)=a.
\]
The three arguments are distinct nonzero algebraic numbers, so this relation
contradicts \cref{thm:input-independence}.

Splitting the defining integral at (2) and (3), and using
(1/t>1/2) on (1<t<2) and (1/t>1/3) on (2<t<3), gives
\[
 E_1(1)>\frac1{2e}+\frac{e^{-2}-e^{-3}}3>\frac15.
\]
Thus \(a=1/m\) is admissible for \(m\ge5\).  Finally, writing \(q=1+h\),
direct differentiation at one gives
\[
 \mathscr R(1+h)=e^{-1}(h^2-h^3+O(h^4)).
\]
Series inversion at \(\mathscr R=1/m\) proves
\eqref{eq:rm-asymptotic}.
\end{proof}

\subsection{A logarithmic lattice and density-one transcendence}

For \(r\ge1\), put
\begin{equation}\label{eq:cr}
 c_r=\Cin(r)-\log r=\gamma-\Ci(r).
\end{equation}
These levels approach \(\gamma\) only polynomially, but the loss caused by
the logarithms can be counted by primes.

\begin{lemma}\label{lem:cr-positive}
One has \(c_r>0\) for every integer \(r\ge1\).
\end{lemma}

\begin{proof}
For \(x>0\), set \(c(x)=\gamma-\Ci(x)\).  Since
\(c'(x)=-\cos x/x\), its local minima occur at
\(x_k=\pi/2+2\pi k\).  At the first minimum,
\[
 \Cin(\pi/2)
 \ge \int_0^{\pi/2}\left(\frac t2-\frac{t^3}{24}\right)dt
 =\frac{\pi^2}{16}-\frac{\pi^4}{1536}
 >\log(\pi/2),
\]
using \(1-\cos t\ge t^2/2-t^4/24\) on this interval.  The last strict
inequality follows, for example, from
\(333/106<\pi<355/113\) and
\(\log(1+u)\le u-u^2/2+u^3/3\) for \(0<u<1\).
Thus
\(c(x_0)>0\).  For \(k\ge1\), integration by parts gives
\(\Ci(x_k)\le2/x_k<1/2<\gamma\).  Here we used
\(x_k\ge5\pi/2\) and the elementary bound \(\gamma>1/2\).
Hence all local minima are positive, and so is
\(c(x)\).  In particular \(c_r>0\).
\end{proof}

For each \(r\), let
\begin{equation}\label{eq:Ur}
 U_r=T_2(c_r)=\frac1{c_r^2}-\frac1{2c_r}.
\end{equation}

\begin{theorem}[Prime-defect lattice theorem]
\label{thm:prime-defect}
Let \(\pi(M)\) be the number of primes not exceeding \(M\).  Then
\begin{equation}\label{eq:prime-defect-trdeg}
 \trdeg_KK(c_1,\ldots,c_M)
 =\trdeg_KK(U_1,\ldots,U_M)
 \ge M-\pi(M).
\end{equation}
Consequently at most \(\pi(M)\) of the first \(M\) members of each family
\(\{c_r\}\) and \(\{U_r\}\) are algebraic over \(K\).  Both families are
\(K\)-transcendental on a set of natural density one.
\end{theorem}

\begin{proof}
The companion result \cite[Cor.~5.5]{KEIO2026I} makes
\(\Cin(1),\ldots,\Cin(M)\) algebraically independent over \(K\).
Equivalently, this follows from
\(\Cin(r)=\{\Ein(ir)+\Ein(-ir)\}/2\): the pairs
\(\{ir,-ir\}\) are disjoint, so these averaging rows have full rank.
Since
\[
 \log r=\sum_{p\le M}v_p(r)\log p,
\]
we have
\[
 K(\Cin(1),\ldots,\Cin(M))
 \subset K(c_1,\ldots,c_M,\log p:p\le M).
\]
Therefore
\[
 M\le \trdeg_KK(c_1,\ldots,c_M)+\pi(M),
\]
which gives the lower bound for the \(c_r\).  Equation
\[
 2U_rc_r^2+c_r-2=0
\]
shows that \(K(c_1,\ldots,c_M)\) is algebraic over
\(K(U_1,\ldots,U_M)\), while the reverse field inclusion is rational.
Thus the two transcendence degrees agree.  If \(a\) of the first \(M\)
coordinates were algebraic over \(K\), the transcendence degree would be at
most \(M-a\); hence \(a\le\pi(M)\).  The prime number theorem completes the
density statement.
\end{proof}

The trace \(U_r\) also converges to \(S_2\).  Direct subtraction gives
\begin{equation}\label{eq:Ur-S2-exact}
 U_r-S_2=
 \frac{\Ci(r)\{\gamma(4-\gamma)-(2-\gamma)\Ci(r)\}}
 {2\gamma^2(\gamma-\Ci(r))^2},
\end{equation}
and the standard expansion of \(\Ci\) yields
\begin{align}
 U_r-S_2
 &=\frac{4-\gamma}{2\gamma^3}\frac{\sin r}{r}
 -\frac{4-\gamma}{2\gamma^3}\frac{\cos r}{r^2}\notag\\
 &\quad+
 \frac{6-\gamma}{2\gamma^4}\frac{\sin^2r}{r^2}
 +O(r^{-3}).\label{eq:Ur-S2-asymptotic}
\end{align}
Thus this lattice is generally only \(O(r^{-1})\) from the Euler trace.
The theorem does not decide any specified \(c_r,U_r\) with \(r\ge2\);
\(c_1=\Cin(1)\) is already \(K\)-transcendental.

\section{Normalized certificates and cross-index rigidity}
\label{sec:certificates}

The moving towers concern limits.  We now return to finite witnesses and
show that all normalized linear certificates are unique.

Let
\[
 V=\bigoplus_{\alpha\in\Qbar^\times}K^{\mathrm{alg}}[\alpha]
\]
be the free \(K^{\mathrm{alg}}\)-vector space with finite support on the
nonzero algebraic points, and define
\[
 \Phi:V\longrightarrow\C,
 \qquad \Phi([\alpha])=\Ein(\alpha).
\]

\begin{lemma}[The evaluation intersection]
\label{lem:evaluation-intersection}
The map \(\Phi\) is injective and
\begin{equation}\label{eq:Phi-intersection}
 \Phi(V)\cap K^{\mathrm{alg}}=\{0\}.
\end{equation}
\end{lemma}

\begin{proof}
By \cref{thm:input-independence}, every finite collection of the displayed
\(\Ein\)-values is algebraically independent over \(K^{\mathrm{alg}}\), and
hence linearly independent.  Moreover any nonzero linear form in
algebraically independent variables is transcendental over the base field:
extend it to an invertible linear coordinate system.  This proves both
claims.
\end{proof}

\begin{definition}
A finite \(K^{\mathrm{alg}}\)-certificate for \(\gamma\) is a triple
\((v,c,a)\in V\times(K^{\mathrm{alg}})^\times\times K^{\mathrm{alg}}\)
satisfying
\begin{equation}\label{eq:certificate}
 \Phi(v)=c\gamma+a.
\end{equation}
Its normalized data are \((v/c,a/c)\).
\end{definition}

\begin{theorem}[Normalized certificate uniqueness]
\label{thm:certificate-uniqueness}
If \((v_1,c_1,a_1)\) and \((v_2,c_2,a_2)\) are two certificates, then
\begin{equation}\label{eq:normalized-certificate-unique}
 \boxed{\displaystyle
 \frac{v_1}{c_1}=\frac{v_2}{c_2},\qquad
 \frac{a_1}{c_1}=\frac{a_2}{c_2}.}
\end{equation}
\end{theorem}

\begin{proof}
Subtracting the two normalized identities gives
\[
 \Phi\!\left(\frac{v_1}{c_1}-\frac{v_2}{c_2}\right)
 =\frac{a_1}{c_1}-\frac{a_2}{c_2}\in K^{\mathrm{alg}}.
\]
Apply \eqref{eq:Phi-intersection}.
\end{proof}

This is an exclusion theorem, not an existence theorem.  The certificate
set may be empty.  Its force is that whenever two apparently unrelated
algebraicity events would produce certificates, their normalized supports
must coincide exactly.

We apply it to the real Euler witnesses recalled in the introduction.  The
existence, location and within-index algebraic-zero rigidity of \(q_*\),
\(\beta_n\) and \(\lambda_n\) were established in
\cite[Thm.~14.2, Prop.~14.3, and Thm.~14.4]{KEIO2026I}.

\begin{theorem}[Cross-index witness rigidity]
\label{thm:cross-index}
At most one member of
\begin{equation}\label{eq:witness-set}
 \mathcal W=
 \{q_*\}\cup\{\beta_n:n\ge2\}
 \cup\{\lambda_n:n\ge3,\ n\text{ odd}\}
\end{equation}
is algebraic.  Thus precisely one of the following alternatives holds:
\begin{enumerate}[label=\textup{(\roman*)}]
 \item every member of \(\mathcal W\) is transcendental;
 \item exactly one member is algebraic, and in that event \(\gamma\) and
 \(\delta\) are algebraically independent over \(K\).
\end{enumerate}
\end{theorem}

\begin{proof}
If \(q_*\) is algebraic, its normalized certificate vector is
\[
 v_*=[q_*],\qquad \Phi(v_*)=\gamma.
\]
If an algebraic number \(\alpha\) is a zero of \(G_n\), its normalized
certificate vector is
\begin{equation}\label{eq:G-certificate-vector}
 v_{n,\alpha}
 =\frac n{n-1}[\alpha]-\frac1{n-1}[\alpha^n],
 \qquad \Phi(v_{n,\alpha})=\gamma.
\end{equation}
The real-zero classification places \(\alpha\) in \((0,1)\) or, for odd
\(n\), in \((-\infty,-1)\).  Hence \(\alpha\ne\alpha^n\), and
\(v_{n,\alpha}\) has exactly one positive and one negative support
coefficient.  It cannot equal the one-point vector \(v_*\).

If \(v_{n,\alpha}=v_{m,\beta}\), equality of the positive and negative
support coordinates gives
\[
 \frac n{n-1}=\frac m{m-1},
\]
so \(n=m\), and then \(\alpha=\beta\).  Certificate uniqueness therefore
excludes two distinct algebraic members of \eqref{eq:witness-set}.  The
algebraic independence of \(\gamma,\delta\) in the exceptional case is
\cite[Thm.~14.2 and Thm.~14.4]{KEIO2026I}.
\end{proof}

\subsection{A defect-one family of critical cancellation coefficients}

Fix an algebraic number \(q>1\).  For \(N\ge1\), put
\begin{align}
 D_N(q)&=E_1(q^N)-2E_1(q^{N+1})+E_1(q^{N+2}),\label{eq:DNq}\\
 A_N(q)&=(N+2)E_1(q^{N+1})-(N+1)E_1(q^{N+2}),\label{eq:ANq}\\
 u_N(q)&=\frac{A_N(q)}{D_N(q)}.\label{eq:uNq}
\end{align}
Both numerator and denominator are positive: \(D_N(q)>0\) is the strict
second-difference inequality for the strictly convex function
\(t\mapsto E_1(q^t)\), while monotonicity gives
\(A_N(q)>E_1(q^{N+1})>0\).  The number \(u_N(q)\) is the
unique positive coefficient for which the three-point augmentation-one tail
vanishes:
\begin{equation}\label{eq:uN-tail-zero}
 -u_NE_1(q^N)+(N+2+2u_N)E_1(q^{N+1})
 -(N+1+u_N)E_1(q^{N+2})=0.
\end{equation}

\begin{theorem}[One-defect cancellation family]
\label{thm:uN-defect}
For every finite set \(I\subset\N_{\ge1}\),
\begin{equation}\label{eq:uN-trdeg}
 \boxed{\displaystyle
 \trdeg_K K(u_N(q):N\in I)\ge |I|-1.}
\end{equation}
Consequently at most one member of \(\{u_N(q):N\ge1\}\) is algebraic over
\(K\).  If one member is algebraic over \(K\), all the remaining members are
countably algebraically independent over \(K\), and \(\gamma,\delta\) are
algebraically independent over \(K\).  If \(\gamma\) is algebraic over
\(K\), then the whole family \(\{u_N(q):N\ge1\}\) is countably
algebraically independent over \(K\).

Finally,
\begin{equation}\label{eq:uN-asymptotic}
 u_N(q)\sim\frac{N+2}{q}
 \exp\!\bigl(-(q-1)q^N\bigr).
\end{equation}
\end{theorem}

\begin{proof}
Let \(X_j=\Ein(q^j)\).  The connection formula and
\eqref{eq:uN-tail-zero} give
\begin{equation}\label{eq:uN-certificate}
 -u_NX_N+(N+2+2u_N)X_{N+1}
 -(N+1+u_N)X_{N+2}=\gamma.
\end{equation}
For finite \(I\), let
\(U=\bigcup_{N\in I}\{N,N+1,N+2\}\).  Since \(u_N>0\), the relations
\eqref{eq:uN-certificate}, used in decreasing order of \(N\), solve every
\(X_N\) with \(N\in I\) in terms of \(u_N\), \(\gamma\), and the
coordinates \(X_j\) with \(j\in U\setminus I\).  Hence
\[
 K(X_j:j\in U)
 \subset
 K(u_N:N\in I,\gamma,X_j:j\in U\setminus I).
\]
The left-hand coordinates are algebraically independent over \(K\), so
\[
 |U|\le
 \trdeg_KK(u_N:N\in I)+1+|U|-|I|,
\]
which proves \eqref{eq:uN-trdeg}.  If one \(u_{N_0}\) is algebraic over
\(K\), apply the bound to \(I\cup\{N_0\}\) to obtain full independence of
every finite set with \(N_0\) omitted.

In that exceptional case \eqref{eq:uN-certificate} expresses \(\gamma\) as
a nonconstant \(K^{\rm alg}\)-linear form in three algebraically independent
\(\Ein\)-coordinates.  Together with
\(\delta/e=\Ein(1)-\gamma\), where \(1\) is a fourth distinct algebraic
point, two linearly independent coordinate forms are obtained.  Thus
\(\gamma,\delta\) are algebraically independent over \(K\).  If instead
\(\gamma\in K^{\rm alg}\), the extra ``\(+1\)'' in the preceding
transcendence-degree estimate disappears, proving full independence of the
\(u_N\).

The asymptotic follows from \eqref{eq:E1-asymptotic}: the first terms in
\eqref{eq:DNq} and \eqref{eq:ANq} dominate and
\[
 \frac{E_1(q^{N+1})}{E_1(q^N)}
 \sim q^{-1}\exp\!\bigl(-(q-1)q^N\bigr).
\]
\end{proof}

\section{Relative exclusion and finite trace descent}
\label{sec:relative}

Put
\[
 A=\Ein(1),\qquad B=\Ein(-1),\qquad
 C=\Cin(1)=\frac{\Ein(i)+\Ein(-i)}2.
\]
The numbers (A,B,C) are algebraically independent over \(K\), and
\begin{align}
 E_1(1)&=A-\gamma,&
 \Ei(1)&=\gamma-B,\label{eq:six-identities}\\
 \Ci(1)&=\gamma-C,&
 \Chi(1)&=\gamma-\frac{A+B}{2},&
 \delta&=e(A-\gamma).\notag
\end{align}

\begin{theorem}[Five relative certificate classes]
\label{thm:five-classes}
Among the following five assertions, at most one can hold:
\begin{align*}
 \mathcal C_0&:\ \gamma\in K^{\mathrm{alg}},\\
 \mathcal C_1&:\ E_1(1)\in K^{\mathrm{alg}}
 \quad\Longleftrightarrow\quad \delta\in K^{\mathrm{alg}},\\
 \mathcal C_2&:\ \Ei(1)\in K^{\mathrm{alg}},\\
 \mathcal C_3&:\ \Ci(1)\in K^{\mathrm{alg}},\\
 \mathcal C_4&:\ \Chi(1)\in K^{\mathrm{alg}}.
\end{align*}
Consequently at least four members of
\[
 \gamma,\quad\delta,\quad E_1(1),\quad\Ei(1),\quad\Ci(1),\quad\Chi(1)
\]
are transcendental over \(K\).  If the possible exceptional class is not
\(\mathcal C_1\), then at least five are transcendental over \(K\).
\end{theorem}

\begin{proof}
The equivalence in \(\mathcal C_1\) follows from
\(\delta=eE_1(1)\) and \(e\in K^\times\).  The five assertions produce the
following normalized certificate vectors, respectively:
\[
 0,\qquad [1],\qquad[-1],\qquad
 \frac{[i]+[-i]}2,\qquad
 \frac{[1]+[-1]}2.
\]
For example, if \(E_1(1)\in K^{\mathrm{alg}}\), then
\(\Phi([1])=\gamma+E_1(1)\); if \(\Ci(1)\in K^{\mathrm{alg}}\), then
\(\Phi(([i]+[-i])/2)=\gamma-\Ci(1)\).  These five vectors are distinct, so
\cref{thm:certificate-uniqueness} permits at most one assertion.  Counting
the sizes (1,2,1,1,1) of the five classes gives the final statement.
\end{proof}

\begin{remark}[The necessary correction]
The theorem does \emph{not} say that at most one of the six constants is
algebraic over \(K\).  The pair \(E_1(1),\delta\) forms one relative class:
they are simultaneously algebraic or transcendental over \(K\).  Over
\(\Qbar\), the companion paper excludes their simultaneous algebraicity by
the additional identity \(e=\delta/E_1(1)\) and the transcendence of \(e\).
\end{remark}

\begin{corollary}[Witnesses exclude every relative class]
\label{cor:witness-all-K-trans}
If one member of the witness set \(\mathcal W\) in
\eqref{eq:witness-set} is algebraic, then all six constants displayed in
\cref{thm:five-classes} are transcendental over \(K\).
\end{corollary}

\begin{proof}
The normalized vector of an algebraic witness is either \([q_*]\) or a
two-point vector with one positive and one negative coefficient as in
\eqref{eq:G-certificate-vector}.  Its support and sign pattern differ from
each of the five vectors in the proof of \cref{thm:five-classes}.  Certificate
uniqueness therefore excludes all five relative assertions.
\end{proof}

We next separate finite field membership from countable analytic limits.
Let
\begin{equation}\label{eq:G-field}
 \calG=K\bigl(\Ein(\alpha):\alpha\in\Qbar^\times\bigr).
\end{equation}
Every element of \(\calG\) uses only finitely many algebraic-point values,
although the generating family is countable.

\begin{proposition}[Relative trace descent]
\label{prop:relative-trace-descent}
For \(m\ge2\),
\[
 S_m=P_m(\gamma^{-1}),\qquad
 \calG(S_m)\subset\calG(\gamma),\qquad
 [\calG(\gamma):\calG(S_m)]\le m.
\]
Consequently
\begin{equation}\label{eq:relative-trdeg}
 \trdeg_{\calG}\calG(S_m)=\trdeg_{\calG}\calG(\gamma).
\end{equation}
Moreover
\begin{equation}\label{eq:G-S2-S3}
 \boxed{\calG(S_2,S_3)=\calG(S_2,S_4)=\calG(\gamma)},
\end{equation}
and hence
\begin{equation}\label{eq:joint-gamma-free}
 S_2,S_3\in\calG
 \quad\Longleftrightarrow\quad
 S_2,S_4\in\calG
 \quad\Longleftrightarrow\quad
 \gamma\in\calG.
\end{equation}
Individually,
\[
 S_2\in\calG\Longrightarrow[\calG(\gamma):\calG]\le2,
 \qquad
 S_3\in\calG\Longrightarrow[\calG(\gamma):\calG]\le3,
 \qquad
 S_4\in\calG\Longrightarrow[\calG(\gamma):\calG]\le4.
\]
\end{proposition}

\begin{proof}
The polynomial identity comes from \cref{thm:universal-trace} at the Euler
level.  Thus \(\gamma^{-1}\) is algebraic of degree at most \(m\) over
\(\calG(S_m)\), proving the first assertions and
\eqref{eq:relative-trdeg}.  Equations \eqref{eq:gamma-recovery} and
\eqref{eq:gamma-even-recovery} prove \eqref{eq:G-S2-S3}; the remaining
statements follow.
\end{proof}

The reverse of the first individual implication for \(S_2\) is false without
an additional norm--trace condition.

\begin{theorem}[Exact quadratic descent criterion]
\label{thm:quadratic-descent}
One has \(S_2\in\calG\) if and only if either \(\gamma\in\calG\), or
\[
 [\calG(\gamma):\calG]=2
 \quad\text{and}\quad
 \Norm_{\calG(\gamma)/\calG}(\gamma)
 =2\Tr_{\calG(\gamma)/\calG}(\gamma).
\]
In the quadratic case the conjugate is
\begin{equation}\label{eq:gamma-conjugate}
 \gamma'=-\frac{2\gamma}{2-\gamma}.
\end{equation}
\end{theorem}

\begin{proof}
Let \(f(x)=x^{-2}-(2x)^{-1}\).  If \(\gamma\notin\calG\) and
\(f(\gamma)=S_2\in\calG\), then \cref{prop:relative-trace-descent} gives a
quadratic extension.  If \(\gamma'\) is its nontrivial conjugate, then
\[
 f(\gamma)-f(\gamma')
 =\frac{(\gamma'-\gamma)
 \{2(\gamma+\gamma')-\gamma\gamma'\}}
 {2\gamma^2{\gamma'}^2}.
\]
Thus invariance of \(f(\gamma)\) is equivalent to
\(\gamma\gamma'=2(\gamma+\gamma')\), which is exactly the displayed
norm--trace condition and rearranges to \eqref{eq:gamma-conjugate}.
Conversely that condition makes \(f(\gamma)\) fixed by the quadratic Galois
involution and hence an element of \(\calG\).
\end{proof}

For instance, degree two alone is insufficient: with base field \(\Q\) and
\(x=\sqrt2\), one has
\(f(x)=1/2-\sqrt2/4\notin\Q\).
